% ============================================================
%  يُترجَم بـ XeLaTeX  (في Overleaf: Menu → Compiler → XeLaTeX)
%  الخط: Amiri (متوفر في Overleaf، أو حمّله من github.com/aliftype/amiri)
% ============================================================
\documentclass[12pt,a4paper]{article}
\usepackage[a4paper,margin=1.7cm,top=2.3cm,bottom=2.9cm,headheight=15pt,footskip=1.5cm]{geometry}
\usepackage{amsmath,amssymb}
\usepackage{icomma}
\usepackage[english,bidi=default,layout=lists]{babel}
\babelprovide[import,main]{arabic}
\babelfont{rm}[Scale=1.08]{Amiri}
\babelfont{sf}[Scale=1.08]{Amiri}
\usepackage[most]{tcolorbox}
\usepackage{tikz}
\usetikzlibrary{arrows.meta,calc,shadings}
\usepackage{array}
\usepackage{colortbl}
\usepackage{fancyhdr}
\usepackage{titlesec}
\usepackage{enumitem}
\usepackage{collcell}
\usepackage{needspace}
\usepackage{etoolbox}
\newcolumntype{C}{>{\collectcell\ar}c<{\endcollectcell}}
\pretocmd{\section}{\Needspace{15\baselineskip}}{}{}
\pretocmd{\subsection}{\Needspace{11\baselineskip}}{}{}

\definecolor{cDef}{HTML}{1F4E9C}
\definecolor{cThm}{HTML}{B3261E}
\definecolor{cRem}{HTML}{2E7D32}
\definecolor{cEx}{HTML}{C76B00}
\definecolor{cMeth}{HTML}{6A1B9A}
\definecolor{cSol}{HTML}{00695C}
\definecolor{cSec}{HTML}{1F3A68}

\titleformat{\section}{\Large\bfseries\color{cSec}}{\thesection.}{0.6em}{}[{\color{cSec!60}\titlerule[1pt]}]
\titleformat{\subsection}{\large\bfseries\color{cThm!85!black}}{\thesubsection.}{0.5em}{}
\titlespacing*{\section}{0pt}{14pt}{8pt}
\titlespacing*{\subsection}{0pt}{10pt}{5pt}
\setlength{\parindent}{0pt}
\setlength{\parskip}{4pt}
\setlist{itemsep=2pt,topsep=3pt}

\newtcolorbox{arbox}[3][]{enhanced,
  colback=#2!5,colframe=#2!80!black,arc=2mm,boxrule=0.8pt,
  left=6pt,right=6pt,top=8pt,bottom=5pt,
  fonttitle=\bfseries,coltitle=white,
  title={\foreignlanguage{arabic}{#3}},
  attach boxed title to top right={xshift=-5mm,yshift=-3mm},
  boxed title style={colback=#2!80!black,arc=1.5mm,boxrule=0pt},
  before upper={\selectlanguage{arabic}},
  before skip=8pt,after skip=8pt,breakable,#1}
\newenvironment{BOX}[2]{\begin{otherlanguage}{english}\begin{arbox}{#1}{#2}}{\end{arbox}\end{otherlanguage}}
\newenvironment{tarif}[1][تعريف]{\begin{BOX}{cDef}{#1}}{\end{BOX}}
\newenvironment{khassia}[1][خاصية]{\begin{BOX}{cThm}{#1}}{\end{BOX}}
\newenvironment{molahada}[1][ملاحظة]{\begin{BOX}{cRem}{#1}}{\end{BOX}}
\newenvironment{mithal}[1][مثال]{\begin{BOX}{cEx}{#1}}{\end{BOX}}
\newenvironment{tariqa}[1][طريقة]{\begin{BOX}{cMeth}{#1}}{\end{BOX}}

\newcounter{exo}
\newenvironment{tamrin}[1]{\stepcounter{exo}\begin{BOX}{cSec}{التمرين \theexo: #1}}{\end{BOX}}
\newenvironment{hal}{\Needspace{12\baselineskip}\begin{otherlanguage}{english}\begin{arbox}[breakable]{cSol}{حل التمرين \theexo}}{\end{arbox}\end{otherlanguage}}

\newenvironment{LTR}{\begin{otherlanguage}{english}}{\end{otherlanguage}}
\newcommand{\ar}[1]{\foreignlanguage{arabic}{#1}}

\renewcommand{\dots}{\mathinner{\ldotp\ldotp\ldotp}}
\renewcommand{\ldots}{\mathinner{\ldotp\ldotp\ldotp}}
\renewcommand{\cdots}{\mathinner{\cdotp\cdotp\cdotp}}
\newcommand{\N}{\mathbb{N}}
\newcommand{\Z}{\mathbb{Z}}
\newcommand{\D}{\mathbb{D}}
\newcommand{\Q}{\mathbb{Q}}
\newcommand{\R}{\mathbb{R}}
\newcommand{\pgcd}{\operatorname{PGCD}}
\newcommand{\ppcm}{\operatorname{PPCM}}

% \NLseg{a}{b}{[}{]} ; \NLright{a}{[} ; \NLleft{b}{]}  (interval on a number line)
\newcommand{\NLseg}[4]{\begin{LTR}\begin{tikzpicture}[baseline=-0.6ex]
 \draw[-{Stealth[length=2mm]},gray!70](0,0)--(3.4,0);
 \draw[line width=2pt,cThm!80!black](0.7,0)--(2.7,0);
 \node[cThm!80!black,font=\LARGE] at (0.7,0){$#3$};
 \node[cThm!80!black,font=\LARGE] at (2.7,0){$#4$};
 \node[below=3pt,font=\footnotesize] at (0.7,0){$#1$};
 \node[below=3pt,font=\footnotesize] at (2.7,0){$#2$};
\end{tikzpicture}\end{LTR}}
\newcommand{\NLright}[2]{\begin{LTR}\begin{tikzpicture}[baseline=-0.6ex]
 \draw[-{Stealth[length=2mm]},gray!70](0,0)--(3.4,0);
 \draw[line width=2pt,cThm!80!black](1.3,0)--(3.3,0);
 \node[cThm!80!black,font=\LARGE] at (1.3,0){$#2$};
 \node[below=3pt,font=\footnotesize] at (1.3,0){$#1$};
\end{tikzpicture}\end{LTR}}
\newcommand{\NLleft}[2]{\begin{LTR}\begin{tikzpicture}[baseline=-0.6ex]
 \draw[-{Stealth[length=2mm]},gray!70](0,0)--(3.4,0);
 \draw[line width=2pt,cThm!80!black](0.05,0)--(2.1,0);
 \node[cThm!80!black,font=\LARGE] at (2.1,0){$#2$};
 \node[below=3pt,font=\footnotesize] at (2.1,0){$#1$};
\end{tikzpicture}\end{LTR}}

% ---------- social-media icons (pure TikZ) ----------
\definecolor{cFB}{HTML}{1877F2}
\definecolor{cTG}{HTML}{229ED9}
\definecolor{cIG1}{HTML}{FEDA75}
\definecolor{cIG2}{HTML}{D62976}
\definecolor{cIG3}{HTML}{4F5BD5}
\newcommand{\fbicon}{\begin{LTR}\begin{tikzpicture}[baseline=-0.75ex,x=1em,y=1em]
 \fill[cFB] (0,0) circle (0.55);
 \node[white,inner sep=0,font=\usefont{TU}{lmss}{bx}{n}\small] at (0.04,-0.02) {f};
\end{tikzpicture}\end{LTR}}
\newcommand{\igicon}{\begin{LTR}\begin{tikzpicture}[baseline=-0.75ex,x=1em,y=1em]
 \shade[lower left=cIG1,lower right=cIG2,upper left=cIG2,upper right=cIG3,rounded corners=0.28em] (-0.55,-0.55) rectangle (0.55,0.55);
 \draw[white,line width=0.09em,rounded corners=0.18em] (-0.36,-0.36) rectangle (0.36,0.36);
 \draw[white,line width=0.09em] (0,0) circle (0.17);
 \fill[white] (0.22,0.22) circle (0.05);
\end{tikzpicture}\end{LTR}}
\newcommand{\tgicon}{\begin{LTR}\begin{tikzpicture}[baseline=-0.75ex,x=1em,y=1em]
 \fill[cTG] (0,0) circle (0.55);
 \fill[white] (-0.33,0.0) -- (0.32,0.27) -- (0.2,-0.3) -- (0.0,-0.15) -- (-0.08,-0.28) -- (-0.08,-0.08) -- cycle;
 \fill[cTG!70!black] (-0.08,-0.08) -- (0.25,0.2) -- (0.0,-0.15) -- (-0.08,-0.28) -- cycle;
\end{tikzpicture}\end{LTR}}
% page number: black box, white two-digit number, at the right of the social row
\newcommand{\PageBox}{\begin{LTR}\tikz[baseline=(p.base)]\node[fill=black,text=white,inner sep=2.5pt,minimum width=1.6em,minimum height=1.3em,font=\usefont{TU}{lmss}{bx}{n}\small](p){\ifnum\value{page}<10 0\fi\arabic{page}};\end{LTR}}
% visual left-to-right order: Telegram | Instagram | Facebook | page box
\newcommand{\SocialRow}{\begin{otherlanguage}{english}%
 \ar{\textbf{الأستاذ سوايسية هشام للرياضيات}}\ \tgicon\hfill
 \textbf{prof\_hicham\_math}\ \igicon\hfill
 \ar{\textbf{الأستاذ سوايسية هشام للرياضيات}}\ \fbicon\hspace{2.2em}\PageBox\end{otherlanguage}}

\pagestyle{fancy}
\fancyhf{}
\renewcommand{\headrulewidth}{0.6pt}
\renewcommand{\footrulewidth}{0.4pt}
% \SetHeader{right head: document title}{left head: level -- stream}
\newcommand{\SetHeader}[2]{%
 \fancyhead[C]{\parbox{\textwidth}{\ar{#1\hfill #2}}}%
 \fancyfoot[C]{\parbox{\textwidth}{\small\SocialRow}}}
\usepackage{tkz-tab}
\usepackage{multicol}
% variation / sign table (LTR)
\newenvironment{TAB}{\par\begin{center}\begin{LTR}\begin{tikzpicture}[scale=0.95,every node/.style={font=\small}]}{\end{tikzpicture}\end{LTR}\end{center}}
% grid + axes, clipped: \begin{GR}{xmin}{xmax}{ymin}{ymax}{unit cm}
\newenvironment{GR}[5]{\par\begin{center}\begin{LTR}\begin{tikzpicture}[x=#5cm,y=#5cm,font=\footnotesize]
 \draw[gray!25,very thin] (#1,#3) grid (#2,#4);
 \draw[-{Stealth[length=2mm]}] (#1,0)--(#2+0.3,0) node[below left]{$x$};
 \draw[-{Stealth[length=2mm]}] (0,#3)--(0,#4+0.3) node[below left]{$y$};
 \node[below left] at (0,0){$O$};
 \begin{scope}\clip (#1,#3) rectangle (#2,#4);}
{\end{scope}\end{tikzpicture}\end{LTR}\end{center}}
\newcommand{\Ticks}[2]{\foreach \t in {#1} \node[below,font=\scriptsize] at (\t,0){$\t$}; \foreach \t in {#2} \node[left,font=\scriptsize] at (0,\t){$\t$};}
\newcommand{\bx}[1]{\ensuremath{\colorbox{cDef!12}{$\displaystyle #1$}}}
\newcommand{\TwoCol}[4]{\par\begin{LTR}\noindent\begin{minipage}[t]{#1\textwidth}\selectlanguage{arabic}#2\end{minipage}\hfill\begin{minipage}[t]{#3\textwidth}\selectlanguage{arabic}#4\end{minipage}\end{LTR}\par}
\newcommand{\RL}[4]{\TwoCol{#3}{#4}{#1}{#2}}  % \RL{w1}{right content}{w2}{left content}
\newcommand{\ok}{$\textcolor{cRem}{\checkmark}$}
\newcommand{\limn}{\displaystyle\lim_{n\to+\infty}}
% sub-questions, each on its own line
\newenvironment{sq}{\begin{itemize}[leftmargin=2.4em,itemsep=3pt]}{\end{itemize}}
\newcommand{\qa}{\item[(أ)]}\newcommand{\qb}{\item[(ب)]}\newcommand{\qc}{\item[(جـ)]}\newcommand{\qd}{\item[(د)]}\newcommand{\qe}{\item[(هـ)]}
\newcommand{\pa}{\par\smallskip\textbf{(أ)} }\newcommand{\pb}{\par\smallskip\textbf{(ب)} }\newcommand{\pc}{\par\smallskip\textbf{(جـ)} }\newcommand{\pd}{\par\smallskip\textbf{(د)} }\newcommand{\pe}{\par\smallskip\textbf{(هـ)} }
\newcommand{\QQ}[1]{\par\smallskip\textbf{#1.} }
\newenvironment{tmr}{\stepcounter{exo}\begin{BOX}{cSec}{تمرين \theexo}}{\end{BOX}}
\newcommand{\ABC}[3]{\\[2pt]\begin{tabular}{p{0.28\textwidth}p{0.28\textwidth}p{0.28\textwidth}}(جـ) \ #3 & (ب) \ #2 & (أ) \ #1\end{tabular}}
\pretocmd{\section}{\Needspace{24\baselineskip}}{}{}
\newtcbox{\bxtbox}{on line,colback=cDef!12,colframe=cDef!12,boxrule=0pt,arc=1pt,boxsep=0pt,left=2pt,right=2pt,top=1pt,bottom=1pt}
\newcommand{\bxt}[1]{\foreignlanguage{english}{\bxtbox{\ar{#1}}}}
\SetHeader{\small\bfseries\color{cSec}تمارين البكالوريا: الأعداد المركبة والتحويلات النقطية}{\small\color{cSec}السنة الثالثة ثانوي -- الشعب العلمية}
\begin{document}
\begin{LTR}
\begin{tcolorbox}[enhanced,colback=cSec,colframe=cSec,arc=3mm,halign=center,boxrule=0pt,top=12pt,bottom=12pt,before upper={\selectlanguage{arabic}}]
{\color{white}\fontsize{24}{30}\selectfont\bfseries تمارين البكالوريا في الأعداد المركبة والتحويلات النقطية}\\[8pt]
{\color{white}\large -- بدون حلول (للتدرّب) -- حتى دورة 2026}\\[5pt]
{\color{white}\large علوم تجريبية -- تقني رياضي -- رياضيات}
\end{tcolorbox}
\end{LTR}
\section{تمارين البكالوريا -- شعبة علوم تجريبية}
% ===================== التمرين 1 =====================
\Needspace{12\baselineskip}
\begin{tamrin}{بكالوريا 2026 -- الموضوع الأول -- علوم تجريبية}
\textbf{I)} حل في المجموعة $\mathbb{C}$ المعادلة ذات المجهول $z$ التالية: $(iz+4)(z^2-4z+16)=0$.

\textbf{II)} المستوي المركب منسوب إلى المعلم المتعامد والمتجانس $(O;\vec u,\vec v)$، و$A$، $B$، $C$ و$D$ نقط من المستوي لاحقاتها على الترتيب $z_A=4$، $z_B=iz_A$، $z_C=2-2i\sqrt3$ و$z_D=-iz_C$.
\QQ{1}
\begin{sq}
\qa أكتب كلا من $z_B$، $z_C$ و$z_D$ على الشكل المثلثي.
\qb علم النقط $A$، $B$، $C$ و$D$ ثم تحقق أنها تنتمي إلى نفس الدائرة التي يطلب تعيين مركزها ونصف قطرها.
\end{sq}
\QQ{2}
\begin{sq}
\qa أبرز أن: $z_D-z_B=\sqrt3\,(z_C-z_A)$ ثم استنتج أن المستقيمين $(BD)$ و$(AC)$ متوازيان.
\qb بين أن: $|z_B-z_A|=|z_D-z_C|$.
\qc استنتج طبيعة الرباعي $ABDC$ ثم عين لاحقة النقطة $G$ مركز ثقله.
\end{sq}
\QQ{3} تحقق أن: $z_D-z_A=i\,(z_B-z_C)$ ثم عين الأعداد الطبيعية $n$ حتى يكون $\left(\dfrac{z_D-z_A}{z_B-z_C}\right)^{n}$ عددا حقيقيا.
\end{tamrin}

% ===================== التمرين 2 =====================
\Needspace{12\baselineskip}
\begin{tamrin}{بكالوريا 2026 -- الموضوع الثاني -- علوم تجريبية}
\QQ{1} نعتبر في المجموعة $\mathbb{C}$ المعادلة ذات المجهول $z$ الآتية: $z^2+\alpha z+4=0 \ \ldots\ (E)$ حيث $\alpha$ عدد حقيقي.
\begin{sq}
\qa أحسب $\left(-1+i\sqrt3\right)^2$ ثم استنتج الجذرين التربيعيين للعدد $-2-2i\sqrt3$.
\qb عين $\alpha$ حتى يكون العدد المركب $-1+i\sqrt3$ حلا للمعادلة $(E)$ ثم استنتج حلها الآخر.
\end{sq}
\QQ{2} نعتبر الأعداد المركبة $z_1$، $z_2$ و$z_3$ حيث: $z_1=-1+i\sqrt3$، $z_2=-1-i\sqrt3$ و$z_3=\sqrt2\,(1-i)$.
\begin{sq}
\qa أكتب كلا من $z_1$، $z_2$ و$z_3$ على الشكل المثلثي.
\qb استنتج الشكل المثلثي للعدد المركب $L$ حيث: $L=\dfrac{z_1\times z_3}{z_2}$.
\qc أكتب $L$ على الشكل الجبري ثم استنتج القيمة المضبوطة لكل من $\cos\dfrac{13\pi}{12}$ و$\sin\dfrac{13\pi}{12}$.
\end{sq}
\end{tamrin}

% ===================== التمرين 3 =====================
\Needspace{12\baselineskip}
\begin{tamrin}{بكالوريا 2025 -- الموضوع الأول -- علوم تجريبية}
\QQ{1} حل في مجموعة الأعداد المركبة $\mathbb{C}$ المعادلة ذات المجهول $z$ التالية: $(iz+2)\left(z^2+2\sqrt3\,z+4\right)=0$.
\QQ{2} المستوي المركب منسوب إلى المعلم المتعامد والمتجانس $(O;\vec u,\vec v)$. $A$، $B$ و$C$ نقط من المستوي لاحقاتها على الترتيب: $z_A=2i$، $z_B=-\sqrt3+i$ و$z_C=\overline{z_B}$.
\begin{sq}
\qa أكتب كلا من $z_A$، $z_B$ و$z_C$ على الشكل المثلثي.
\qb استنتج أن النقط $A$، $B$ و$C$ تنتمي إلى نفس الدائرة التي يطلب تعيين مركزها ونصف قطرها.
\qc حدد طبيعة المثلث $ABC$ ثم عين لاحقة مركز ثقله.
\end{sq}
\QQ{3} $z$ عدد مركب حيث $z=(\cos\theta+i\sin\theta)\,z_A$ و$\theta$ عدد حقيقي مع $0<\theta<\dfrac{\pi}{2}$.

عين قيمة $\theta$ بحيث يكون $\dfrac{5\pi}{6}$ عمدة للعدد $z$.
\end{tamrin}

% ===================== التمرين 4 =====================
\Needspace{12\baselineskip}
\begin{tamrin}{بكالوريا 2025 -- الموضوع الثاني -- علوم تجريبية}
\QQ{1} أنشر $(1+i)^2$ ثم حل في مجموعة الأعداد المركبة $\mathbb{C}$ المعادلة ذات المجهول $z$ التالية: $z^2=8i$.
\QQ{2} المستوي المركب منسوب إلى المعلم المتعامد والمتجانس $(O;\vec u,\vec v)$. $A$، $B$ و$C$ نقط من المستوي لاحقاتها على الترتيب: $z_A=2+2i$، $z_B=-z_A$ و$z_C=\overline{z_B}$.
\begin{sq}
\qa أكتب كلا من $z_A$، $z_B$ و$z_C$ على الشكل المثلثي.
\qb استنتج أن النقط $A$، $B$ و$C$ تنتمي إلى نفس الدائرة التي يطلب تعيين مركزها ونصف قطرها.
\end{sq}
\QQ{3} تحقق أن: $z_A-z_C=i\,(z_B-z_C)$ ثم حدد طبيعة المثلث $ABC$.
\QQ{4} عين $z_D$ و$z_E$ لاحقتي النقطتين $D$ و$E$ على الترتيب حتى تكون النقطة $C$ مركز المربع $ABDE$.
\end{tamrin}

% ===================== التمرين 5 =====================
\Needspace{12\baselineskip}
\begin{tamrin}{بكالوريا 2024 -- الموضوع الأول -- علوم تجريبية}
\textbf{I)} حل في مجموعة الأعداد المركبة $\mathbb{C}$ المعادلة ذات المجهول $z$ الآتية:
\[\left(z-1+2\sqrt3\right)\left[z^2-2\left(1-\sqrt3\right)z+5-2\sqrt3\right]=0 .\]
\textbf{II)} في المستوي المنسوب إلى المعلم المتعامد والمتجانس $(O;\vec i,\vec j)$، نعتبر النقط $A$، $B$ و$C$ التي لاحقاتها على الترتيب $z_A$، $z_B$ و$z_C$ حيث: $z_A=1-\sqrt3+i$، $z_B=1-2\sqrt3$ و$z_C=\overline{z_A}$.
\QQ{1} أكتب كلا من $z_A-1$، $z_C-1$ و$z_B$ على الشكل المثلثي.
\QQ{2} جد لاحقة النقطة $D$ مرجح الجملة المثقلة $\{(A;1),(B;-1),(C;1)\}$.
\QQ{3} بين أن الرباعي $ABCD$ معين.
\end{tamrin}

% ===================== التمرين 6 =====================
\Needspace{12\baselineskip}
\begin{tamrin}{بكالوريا 2024 -- الموضوع الثاني -- علوم تجريبية}
عين الاقتراح الصحيح الوحيد من بين الاقتراحات الثلاثة مع التبرير في كل حالة مما يلي:
\QQ{1} $z$ عدد مركب مرافقه $\overline{z}$. مرافق العدد المركب $z+i$ هو:
\ABC{$\overline{z}-i$}{$\overline{z}+i$}{$z-i$}
\QQ{2} العدد المركب $\left(\dfrac{1+i}{1-i}\right)^{2024}$ يساوي:
\ABC{$1$}{$i$}{$-1$}
\QQ{3} $z$ عدد مركب حيث $z=2\left(1+i\sqrt3\right)$.
من أجل كل عدد طبيعي غير معدوم $n$، نضع: $S_n=\ln|z|+\ln|z|^2+\ldots+\ln|z|^n$، لدينا:
\begin{sq}
\qa $S_n=(n+1)^2\ln2$ \qquad\qquad (ب) $S_n=n(n+1)\ln2$
\qc $S_n=2\left(\dfrac{1-(2\ln2)^n}{1-2\ln2}\right)\ln2$
\end{sq}
\QQ{4} $z$ عدد مركب حيث: $z=\sin\dfrac{\pi}{8}+i\cos\dfrac{\pi}{8}$. الشكل المثلثي للعدد المركب $z$ هو:
\ABC{$z=-\cos\dfrac{\pi}{8}+i\sin\dfrac{\pi}{8}$}{$z=\cos\dfrac{\pi}{8}+i\sin\dfrac{\pi}{8}$}{$z=\cos\dfrac{3\pi}{8}+i\sin\dfrac{3\pi}{8}$}
\end{tamrin}

% ===================== التمرين 7 =====================
\Needspace{12\baselineskip}
\begin{tamrin}{بكالوريا 2023 -- الموضوع الثاني -- علوم تجريبية}
عين الاقتراح الصحيح الوحيد من بين الاقتراحات الثلاثة في كل حالة من الحالات التالية مع التبرير:
\QQ{1} حلا المعادلة $8z^2-4z+1=0$ ذات المجهول $z$ في $\mathbb{C}$ هما:
\ABC{$-\frac14+\frac14i$ و$\frac14-\frac14i$}{$-\frac14+\frac14i$ و$-\frac14-\frac14i$}{$\frac14-\frac14i$ و$\frac14+\frac14i$}
\QQ{2} الشكل الجبري للعدد المركب $\dfrac{1+\sqrt3+i}{1-i}$ هو:
\ABC{$\dfrac{\sqrt3}{2}+i\left(\dfrac{2+\sqrt3}{2}\right)$}{$\dfrac{\sqrt3}{2}-i\left(\dfrac{2+\sqrt3}{2}\right)$}{$\dfrac{\sqrt3}{2}+i\left(\dfrac{-2+\sqrt3}{2}\right)$}
\QQ{3} الجذران التربيعيان للعدد المركب $-8+6i$ هما:
\ABC{$1+3i$ و$-1-3i$}{$1+3i$ و$1-3i$}{$3+i$ و$-3-i$}
\QQ{4} الشكل المثلثي للعدد المركب $\dfrac{1+i}{\sqrt3-i}$ هو:
\ABC{$\sqrt2\left(\cos\frac{\pi}{12}+i\sin\frac{\pi}{12}\right)$}{$\frac{\sqrt2}{2}\left(\cos\frac{5\pi}{12}+i\sin\frac{5\pi}{12}\right)$}{$\frac{\sqrt2}{2}\left(\cos\frac{7\pi}{12}+i\sin\frac{7\pi}{12}\right)$}
\end{tamrin}

% ===================== التمرين 8 =====================
\Needspace{12\baselineskip}
\begin{tamrin}{بكالوريا 2019 -- الموضوع الأول -- علوم تجريبية}
\textbf{I)} حل في المجموعة $\mathbb{C}$ المعادلة: $(z-i)\left(z^2-4z+5\right)=0$.

\textbf{II)} نعتبر في المستوي المنسوب إلى المعلم المتعامد والمتجانس $(O;\vec u,\vec v)$ النقط $A$، $B$ و$C$ التي لاحقاتها: $i$، $2-i$ و$2+i$ على الترتيب.
\QQ{1} أكتب العدد المركب $\dfrac{z_C-z_A}{z_C-z_B}$ على الشكل الأسي، ثم استنتج طبيعة المثلث $ABC$.
\QQ{2} من أجل كل عدد مركب $z$ يختلف عن $2+i$ نضع: $f(z)=\dfrac{iz-1-2i}{2z-4-2i}$.
\begin{sq}
\qa عين المجموعة $(E)$ للنقط $M$ من المستوي ذات اللاحقة $z$ والتي تحقق $|f(z)|=\dfrac12$.
\qb بين أن العدد $\left[f(i)\right]^{1440}$ حقيقي موجب.
\end{sq}
\QQ{3} نعتبر الدوران $r$ الذي مركزه النقطة $C$ وزاويته $\dfrac{\pi}{2}$.
\begin{sq}
\qa عين لاحقة النقطة $D$ صورة النقطة $B$ بالدوران $r$، وبين أن النقط $A$، $D$ و$C$ في استقامية.
\qb استنتج أن النقطة $D$ هي صورة النقطة $A$ بتحويل نقطي بسيط يطلب تعيين عناصره المميزة.
\end{sq}
\end{tamrin}

% ===================== التمرين 9 =====================
\Needspace{12\baselineskip}
\begin{tamrin}{بكالوريا 2019 -- الموضوع الثاني -- علوم تجريبية}
المستوي منسوب إلى المعلم المتعامد والمتجانس $(O;\vec u,\vec v)$.
نعتبر النقط $A$، $B$ و$C$ التي لاحقاتها $z_A$، $z_B$ و$z_C$ على الترتيب حيث: $z_A=\sqrt2+i\sqrt6$، $z_B=\overline{z_A}$ و$z_C=-2z_A$.
\QQ{1}
\begin{sq}
\qa أكتب العدد المركب $z_A$ على الشكل الأسي.
\qb أحسب العدد $\left(\dfrac{z_A}{2\sqrt2}\right)^{2019}+\left(\dfrac{z_B}{2\sqrt2}\right)^{2019}$.
\end{sq}
\QQ{2}
\begin{sq}
\qa $T$ الانسحاب الذي يحول $A$ إلى $C$. عين $z_D$ لاحقة النقطة $D$ صورة النقطة $B$ بالانسحاب $T$.
\qb استنتج طبيعة الرباعي $ABDC$.
\end{sq}
\QQ{3} أكتب العدد المركب $z_C-z_A$ على الشكل الأسي.
\QQ{4} جد قيم العدد الطبيعي $n$ التي يكون من أجلها العدد المركب $\left(\dfrac{-6\sqrt2}{z_C-z_A}\right)^n$ عددا حقيقيا.
\QQ{5} لتكن $M$ نقطة كيفية من المستوي لاحقتها $z$ تختلف عن $A$ وتختلف عن $C$.
عين $(E)$ مجموعة النقط $M$ التي يكون من أجلها $\dfrac{z_A-z}{z_C-z}$ عددا حقيقيا موجبا تماما.
\end{tamrin}

% ===================== التمرين 10 =====================
\Needspace{12\baselineskip}
\begin{tamrin}{بكالوريا 2018 -- الموضوع الأول -- علوم تجريبية}
\QQ{1} حل في مجموعة الأعداد المركبة $\mathbb{C}$ المعادلة ذات المجهول $z$ التالية: $z^2-\sqrt3\,z+1=0$.
\QQ{2} المستوي المركب منسوب إلى المعلم المتعامد والمتجانس $(O;\vec u,\vec v)$. $A$، $B$ و$C$ ثلاث نقط من المستوي لواحقها على الترتيب: $z_A=\dfrac12+i\dfrac{\sqrt3}{2}$، $z_B=\dfrac{\sqrt3}{2}+\dfrac12i$ و$z_C=\overline{z_B}$.

أكتب $z_A$ و$z_B$ على الشكل الأسي ثم عين قيم العدد الطبيعي $n$ بحيث يكون: $\left(\dfrac{z_A}{z_B}\right)^n=\dfrac{1+i\sqrt3}{2}$.
\QQ{3}
\begin{sq}
\qa تحقق أن: $\dfrac{z_B}{z_C}=e^{i\frac{\pi}{3}}$ وحدد طبيعة المثلث $OBC$.
\qb استنتج أن $B$ هي صورة $C$ بدوران $r$ يطلب تعيين عناصره المميزة.
\end{sq}
\QQ{4} نسمي $(\gamma)$ مجموعة النقط $M$ ذات اللاحقة $z$ التي تحقق: $|z|=\left|\overline{z}-\dfrac{\sqrt3+i}{2}\right|$.
عين طبيعة المجموعة $(\gamma)$ ثم عين صورتها بالدوران $r$.
\end{tamrin}

% ===================== التمرين 11 =====================
\Needspace{12\baselineskip}
\begin{tamrin}{بكالوريا 2018 -- الموضوع الثاني -- علوم تجريبية}
\textbf{I)} حل في مجموعة الأعداد المركبة $\mathbb{C}$ المعادلة: $\left(\overline{z}-4+i\right)\left(z^2-4z+5\right)=0$ (يرمز $\overline{z}$ لمرافق العدد $z$).

\textbf{II)} في المستوي المركب المنسوب إلى المعلم المتعامد والمتجانس $(O;\vec u,\vec v)$، نعتبر النقط $A$، $B$ و$C$ لواحقها على الترتيب: $z_A=2+i$، $z_B=4+i$ و$z_C=\overline{z_A}$.
\QQ{1} تحقق أن: $\dfrac{z_B-z_A}{z_C-z_A}=i$، ثم عين قيم العدد الطبيعي $n$ بحيث يكون $\left(\dfrac{z_B-z_A}{z_C-z_A}\right)^n$ تخيليا صرفا.
\QQ{2} $D$ نقطة من المستوي لاحقتها $z_D$ حيث:
$\begin{cases}|z_D-z_A|=|z_B-z_A|\\[2pt] \arg\left(\dfrac{z_D-z_A}{z_B-z_A}\right)=\dfrac{\pi}{3}+2k\pi\quad(k\in\mathbb{Z})\end{cases}$

بين أن المثلث $ABD$ متقايس الأضلاع وأحسب $z_D$.
\QQ{3} أحسب $z_G$ لاحقة النقطة $G$ مركز ثقل المثلث $ABD$ ثم عين نسبة وزاوية التشابه المباشر الذي مركزه $A$ ويحول $G$ إلى $D$.
\QQ{4} عين $(\Gamma)$ مجموعة النقط $M$ ذات اللاحقة $z$ ($M$ تختلف عن $C$) بحيث: $\arg\left(\dfrac{z_G-z}{z_C-z}\right)=\pi+2k\pi$ ($k\in\mathbb{Z}$).
\end{tamrin}

% ===================== التمرين 12 =====================
\Needspace{12\baselineskip}
\begin{tamrin}{بكالوريا 2017 -- الدورة الاستثنائية -- الموضوع الأول -- علوم تجريبية}
\textbf{I)} حل في المجموعة $\mathbb{C}$ المعادلة: $(z-2)\left(z^2+2z+4\right)=0$.

\textbf{II)} المستوي المركب منسوب إلى المعلم المتعامد والمتجانس $(O;\vec u,\vec v)$ حيث $\|\vec u\|=2\,\text{cm}$.
نعتبر النقط $A$، $B$ و$C$ التي لاحقاتها: $z_A=2$، $z_B=-1+i\sqrt3$ و$z_C=\overline{z_B}$.
\QQ{1}
\begin{sq}
\qa أكتب $z_B$ على الشكل الأسي، ثم استنتج الشكل الأسي للعدد المركب $z_C$.
\qb عين مركز ونصف قطر الدائرة المحيطة بالمثلث $ABC$، ثم أنشئ النقط $A$، $B$ و$C$.
\end{sq}
\QQ{2} ليكن $S$ التشابه المباشر الذي مركزه النقطة $O$ ونسبته $\dfrac12$ وزاويته $\dfrac{2\pi}{3}$ ولتكن $F$ صورة $A$ بالتحويل $S$.
\begin{sq}
\qa أكتب العبارة المركبة للتشابه $S$ ثم عين لاحقة كل من $A'$، $B'$ و$C'$ صور النقط $A$، $B$ و$C$ على الترتيب بالتشابه $S$، ثم أنشئ في المعلم السابق النقط $A'$، $B'$ و$C'$.
\qb أحسب بالسنتيمتر المربع مساحة المثلث $A'B'C'$.
\end{sq}
\end{tamrin}

% ===================== التمرين 13 =====================
\Needspace{12\baselineskip}
\begin{tamrin}{بكالوريا 2017 -- الدورة الاستثنائية -- الموضوع الثاني -- علوم تجريبية}
المستوي منسوب إلى المعلم المتعامد والمتجانس $(O;\vec u,\vec v)$.
نعتبر النقط $A$، $B$ و$C$ التي لاحقاتها على الترتيب: $z_A=-3-2i$، $z_B=1+i$ و$z_C=4-3i$.
\QQ{1} عين النسبة والزاوية للتشابه $S$ المباشر ذي المركز $A$ والذي يحول النقطة $B$ إلى النقطة $C$.
\QQ{2} أكتب على الشكل الأسي العدد المركب $\dfrac{z_A-z_B}{z_C-z_B}$، ثم استنتج طبيعة المثلث $ABC$.
\QQ{3} نرمز بـ $G$ إلى مركز ثقل المثلث $ABC$ وبـ $I$ إلى منتصف القطعة $[AC]$.
عين كلا من $z_G$ و$z_I$ لاحقتي النقطتين $G$ و$I$، ثم بين أن النقط $B$، $G$ و$I$ في استقامية.
\QQ{4} نعتبر النقطة $D$ نظيرة النقطة $B$ بالنسبة إلى $I$. حدد بدقة طبيعة الرباعي $ABCD$.
\QQ{5} نعتبر $(\Gamma)$ مجموعة النقط $M$ من المستوي التي تحقق: $\left\|\overrightarrow{MA}+\overrightarrow{MC}\right\|=5\sqrt2$.
\begin{sq}
\qa تحقق أن النقطة $C$ تنتمي إلى $(\Gamma)$.
\qb عين طبيعة المجموعة $(\Gamma)$ ثم أنشئها.
\end{sq}
\end{tamrin}

% ===================== التمرين 14 =====================
\Needspace{12\baselineskip}
\begin{tamrin}{بكالوريا 2017 -- الموضوع الأول -- علوم تجريبية}
\textbf{I)} حل في المجموعة $\mathbb{C}$ المعادلة: $(z+2)\left(z^2-4z+8\right)=0$.

\textbf{II)} المستوي المركب منسوب إلى المعلم المتعامد والمتجانس $(O;\vec u,\vec v)$.
نعتبر النقط $A$، $B$ و$C$ التي لاحقاتها: $z_A=2-2i$، $z_B=\overline{z_A}$ و$z_C=-2$.
\QQ{1} أكتب كلا من العددين $z_A$ و$z_B$ على الشكل الأسي.
\QQ{2} عين $z_D$ لاحقة النقطة $D$ حتى تكون النقطة $B$ مركز ثقل المثلث $ACD$.
\QQ{3} $(\Gamma)$ مجموعة النقط $M$ من المستوي ذات اللاحقة $z$ ($M$ تختلف عن $A$ و$B$) بحيث: $\arg\left(\dfrac{z_B-z}{z_A-z}\right)=\dfrac{\pi}{2}$.
تحقق أن مبدأ المعلم $O$ هو نقطة من $(\Gamma)$ ثم عين طبيعة المجموعة $(\Gamma)$ وأنشئها.
\QQ{4} ليكن $h$ التحاكي الذي مركزه النقطة $C$ ونسبته $2$، و$(\Gamma')$ صورة $(\Gamma)$ بالتحاكي $h$.
عين طبيعة المجموعة $(\Gamma')$ مع تحديد عناصرها المميزة.
\end{tamrin}

% ===================== التمرين 15 =====================
\Needspace{12\baselineskip}
\begin{tamrin}{بكالوريا 2017 -- الموضوع الثاني -- علوم تجريبية}
المستوي المركب منسوب إلى المعلم المتعامد والمتجانس $(O;\vec u,\vec v)$.
أجب بصحيح أو خطأ مع التعليل في كل حالة مما يلي:
\QQ{1} مجموعة حلول المعادلة $\left(\dfrac{z+1-i}{z-i}\right)^2=1$ في المجموعة $\mathbb{C}$ هي: $S=\left\{-\dfrac12+i\right\}$.
\QQ{2} من أجل كل عدد مركب $z$، $(z+2)\left(\overline{z}+2\right)=|z+2|^2$.
\QQ{3} من أجل كل عدد طبيعي $n$، $\left(\dfrac12+i\dfrac{\sqrt3}{2}\right)^{3n}=1$.
\QQ{4} $S$ التشابه المباشر الذي مركزه $\Omega$ ذات اللاحقة $1$ ونسبته $3$ وزاويته $\dfrac{\pi}{2}$.
صورة الدائرة $(C)$ ذات المركز $\omega(0;1)$ ونصف القطر $3$ بالتشابه $S$ هي الدائرة $(C')$ ذات المركز $\omega'(-2;-3)$ ونصف القطر $9$.
\QQ{5} من أجل كل عدد حقيقي $\alpha$: إذا كان $z=(\sin\alpha+i\cos\alpha)\times(\cos\alpha-i\sin\alpha)$ فإن: $\arg(z)=\dfrac{\pi}{2}-2\alpha+2k\pi$ حيث $k$ عدد صحيح.
\end{tamrin}

% ===================== التمرين 16 =====================
\Needspace{12\baselineskip}
\begin{tamrin}{بكالوريا 2016 -- الدورة الثانية -- الموضوع الأول -- علوم تجريبية}
\QQ{1} نضع من أجل كل عدد مركب $z$: $P(z)=z^3-24\sqrt3$.
\begin{sq}
\qa تحقق أن: $P\left(2\sqrt3\right)=0$.
\qb جد العددين الحقيقيين $a$ و$b$ بحيث من أجل كل عدد مركب $z$: $P(z)=\left(z-2\sqrt3\right)\left(z^2+az+b\right)$.
\qc حل في مجموعة الأعداد المركبة $\mathbb{C}$ المعادلة $P(z)=0$.
\end{sq}
\QQ{2} المستوي منسوب إلى المعلم المتعامد والمتجانس $(O;\vec u,\vec v)$. $A$، $B$ و$C$ نقط من المستوي لواحقها على الترتيب: $z_A=-\sqrt3+3i$، $z_B=-\sqrt3-3i$ و$z_C=2\sqrt3$.
\begin{sq}
\qa أكتب على الشكل الجبري العدد المركب $\dfrac{z_C-z_A}{z_B-z_A}$.
\qb بين أنه يوجد دوران $R$ مركزه $A$ ويحول النقطة $B$ إلى النقطة $C$، يطلب تعيين زاويته.
\qc استنتج طبيعة المثلث $ABC$.
\qd عين $z_D$ لاحقة النقطة $D$ صورة النقطة $C$ بالانسحاب الذي شعاعه $\overrightarrow{AB}$، ثم حدد بدقة طبيعة الرباعي $ABDC$.
\end{sq}
\QQ{3} عين $(\Gamma)$ مجموعة النقط $M$ من المستوي ذات اللاحقة غير المعدومة $z$ بحيث: $\arg\left(\dfrac{z}{\overline{z}}\right)=2k\pi$ حيث $k\in\mathbb{Z}$.
\end{tamrin}

% ===================== التمرين 17 =====================
\Needspace{12\baselineskip}
\begin{tamrin}{بكالوريا 2016 -- الدورة الثانية -- الموضوع الثاني -- علوم تجريبية}
نعتبر في المجموعة $\mathbb{C}$ المعادلة ذات المجهول $z$ التالية: $2\overline{z}^{\,3}+3\overline{z}^{\,2}-3\overline{z}+5=0\ \ldots\ (E)$.
يشير الرمز $\overline{z}$ إلى مرافق العدد المركب $z$.
\QQ{1}
\begin{sq}
\qa أثبت أن المعادلة $(E)$ تكافئ المعادلة: $\left(2\overline{z}+5\right)\left(\overline{z}^{\,2}-\overline{z}+1\right)=0$.
\qb حل في المجموعة $\mathbb{C}$ المعادلة $(E)$.
\end{sq}
\QQ{2} في المستوي المركب المنسوب إلى المعلم المتعامد والمتجانس $(O;\vec u,\vec v)$ نعتبر النقط $A$، $B$، $C$ و$D$ التي لواحقها على الترتيب: $z_A=\dfrac12-i\dfrac{\sqrt3}{2}$، $z_B=\overline{z_A}$، $z_C=-1$ و$z_D=-\dfrac52$.
\begin{sq}
\qa أكتب كلا من العددين $z_A$ و$z_B$ على الشكل الأسي.
\qb أنشئ النقط $A$، $B$، $C$ و$D$.
\qc أثبت أن: $z_B-z_C=z_B\,(z_A-z_C)$.
\qd استنتج طبيعة المثلث $ABC$.
\end{sq}
\QQ{3} ليكن $S$ التشابه المباشر الذي مركزه $C$ ونسبته $2$ وزاويته $\dfrac{\pi}{3}$، ولتكن $F$ صورة $A$ بالتحويل $S$.
أنشئ النقطة $F$ ثم حدد طبيعة المثلث $AFC$.
\QQ{4} عين طبيعة المجموعة $(\gamma)$ للنقط $M$ من المستوي ذات اللاحقة $z$ حيث $z+1=kz_B$، لما $k$ يتغير في المجموعة $\mathbb{R}_+$.
\end{tamrin}

% ===================== التمرين 18 =====================
\Needspace{12\baselineskip}
\begin{tamrin}{بكالوريا 2016 -- الدورة الأولى -- الموضوع الأول -- علوم تجريبية}
المستوي المركب منسوب إلى المعلم المتعامد والمتجانس $(O;\vec u,\vec v)$. من أجل كل نقطة $M$ من المستوي لاحقتها العدد المركب $z$ ($z\neq1$) نرفق النقطة $M'$ لاحقتها العدد المركب $z'$ حيث: $z'=\dfrac{z-2}{z-1}$.
\QQ{1} حل في $\mathbb{C}$ المعادلة ذات المجهول $z$: $z'=z$.
\QQ{2} النقطتان $A$ و$B$ لاحقتاهما $z_1$ و$z_2$ على الترتيب حيث: $z_1=1-i$، $z_2=\overline{z_1}$.
\begin{sq}
\qa أكتب $\dfrac{z_2}{z_1}$ على الشكل الأسي.
\qb بين أن النقطة $B$ صورة النقطة $A$ بالدوران $R$ الذي مركزه المبدأ $O$، يطلب تعيين زاوية له.
\end{sq}
\QQ{3} نضع: $z'\neq z$. نعتبر النقطتين $C$ و$D$ لاحقتيهما $2$ و$1$ على الترتيب.
عين $(\Gamma)$ مجموعة النقط $M$ حيث $M'$ تنتمي إلى محور التراتيب ثم أنشئه.
\QQ{4} $h$ التحاكي الذي مركزه المبدأ $O$ ونسبته $2$.
\begin{sq}
\qa عين طبيعة التحويل النقطي $S=h\circ R$ وعناصره المميزة.
\qb أكتب العبارة المركبة للتحويل $S$.
\qc عين ثم أنشئ المجموعة $(\Gamma')$ صورة $(\Gamma)$ بالتحويل $S$.
\end{sq}
\end{tamrin}

% ===================== التمرين 19 =====================
\Needspace{12\baselineskip}
\begin{tamrin}{بكالوريا 2016 -- الدورة الأولى -- الموضوع الثاني -- علوم تجريبية}
\QQ{1} حل في مجموعة الأعداد المركبة $\mathbb{C}$ المعادلة ذات المجهول $z$ التالية: $\left(z-\dfrac{\sqrt3}{2}-\dfrac12i\right)\left(z^2+\sqrt3\,z+1\right)=0$.
\QQ{2} المستوي المركب منسوب إلى المعلم المتعامد والمتجانس $(O;\vec u,\vec v)$.
$A$، $B$ و$C$ ثلاث نقط من المستوي لواحقها على الترتيب: $z_A=\dfrac{\sqrt3}{2}+\dfrac12i$، $z_B=-\dfrac{\sqrt3}{2}+\dfrac12i$ و$z_C=\overline{z_B}$.
\begin{sq}
\qa أكتب $z_A$، $z_B$ و$z_C$ على الشكل الأسي.
\qb بين أنه يوجد تشابه مباشر $S$ مركزه $B$ ويحول النقطة $A$ إلى النقطة $C$، يطلب تعيين عناصره المميزة.
\end{sq}
\QQ{3}
\begin{sq}
\qa عين لاحقة النقطة $D$ حتى يكون الرباعي $ABCD$ متوازي أضلاع، ثم حدد بدقة طبيعته.
\qb عين $(E)$ مجموعة النقط $M$ ذات اللاحقة $z$ والتي تحقق: $|z-z_A|=\left|\overline{z}-z_B\right|$.
\qc عين $(\Gamma)$ مجموعة النقط $M$ ذات اللاحقة $z$ والتي تحقق: $z=z_B+\sqrt3\,e^{i\theta}$ عندما $\theta$ يتغير على $\mathbb{R}$، ثم تحقق أن $A\in(\Gamma)$.
\end{sq}
\end{tamrin}

% ====================== التمرين: 2015 م1 ======================
\begin{tamrin}{بكالوريا 2015 -- الموضوع الأول -- علوم تجريبية}
\textbf{I)} عيّن العددين المركّبين $\alpha$ و$\beta$ حيث:
\[\begin{cases}2\alpha-\beta=-3\\[2pt] 2\overline{\alpha}+\overline{\beta}=-3-2i\sqrt3\end{cases}\]
مع $\overline{\alpha}$ مرافق $\alpha$ و$\overline{\beta}$ مرافق $\beta$.

\textbf{II)} المستوي منسوب إلى المعلم المتعامد والمتجانس $(O;\vec u,\vec v)$. $A$، $B$ و$C$ النقط التي لواحقها على الترتيب:
\[z_A=-\frac32+i\frac{\sqrt3}{2},\qquad z_B=\overline{z_A},\qquad z_A=z_C\,e^{i\frac{\pi}{3}}.\]
\textbf{1)}
\begin{sq}
\qa اكتب $z_A$ و$z_C$ على الشكل الأسّي، ثم عيّن قيم العدد الطبيعي $n$ بحيث يكون $\left(\dfrac{z_A}{z_C}\right)^{n}$ حقيقيًا سالبًا.
\qb تحقّق أنّ العدد المركّب $2\left(\dfrac{z_A}{\sqrt3}\right)^{2015}+\left(\dfrac{z_B}{\sqrt3}\right)^{1962}-\left(\dfrac{z_C}{\sqrt3}\right)^{1435}$ حقيقي.
\end{sq}
\textbf{2)} $D$ النقطة ذات اللاحقة $z_D=1+i$.
\begin{sq}
\qa حدّد النسبة والزاوية للتشابه $S$ الذي مركزه $O$ ويحوّل $D$ إلى $A$.
\qb اكتب $\dfrac{z_A}{z_D}$ على الشكل الجبري، ثم استنتج القيمة المضبوطة لكلّ من $\cos\left(\dfrac{7\pi}{12}\right)$ و$\sin\left(\dfrac{7\pi}{12}\right)$.
\end{sq}
\textbf{3)} عيّن مجموعة النقط $M$ ذات اللاحقة $z$ التي تحقّق: $z=k(1+i)\,e^{i\frac{7\pi}{12}}$ حيث $k$ يمسح $\R^{+}$.
\end{tamrin}

% ====================== التمرين: 2015 م2 ======================
\begin{tamrin}{بكالوريا 2015 -- الموضوع الثاني -- علوم تجريبية}
في المستوي المنسوب إلى المعلم المتعامد والمتجانس $(O;\vec u,\vec v)$ نعتبر النقط $A$، $B$ و$C$ التي لواحقها على الترتيب $z_A$، $z_B$ و$z_C$ حيث:
\[z_A=2e^{i\frac{\pi}{6}},\qquad z_B=-\overline{z_A},\qquad z_C=-(z_A+z_B).\]
\textbf{1)}
\begin{sq}
\qa اكتب كلًّا من العددين المركّبين $z_B$ و$z_C$ على الشكل الأسّي.
\qb استنتج أنّ النقط $A$، $B$ و$C$ تنتمي إلى دائرة $(\gamma)$ يُطلب تعيين مركزها ونصف قطرها.
\qc أنشئ الدائرة $(\gamma)$ والنقط $A$، $B$ و$C$.
\end{sq}
\textbf{2)}
\begin{sq}
\qa تحقّق أنّ: $\dfrac{z_B-z_C}{z_B-z_A}=e^{-i\frac{\pi}{3}}$.
\qb استنتج أنّ المثلث $ABC$ متقايس الأضلاع وأنّ النقطة $O$ مركز ثقل هذا المثلث.
\qc عيّن وأنشئ $(E)$ مجموعة النقط $M$ ذات اللاحقة $z$ حيث: $|z|=\left|z-\sqrt3-i\right|$.
\end{sq}
\textbf{3)}
\begin{sq}
\qa عيّن زاوية للدوران $r$ الذي مركزه $O$ ويحوّل $C$ إلى $A$.
\qb أثبت أنّ صورة $(E)$ بالدوران $r$ هي محور القطعة $[OB]$.
\end{sq}
\end{tamrin}

% ====================== التمرين: 2014 م1 ======================
\begin{tamrin}{بكالوريا 2014 -- الموضوع الأول -- علوم تجريبية}
\textbf{1)} حلّ في مجموعة الأعداد المركّبة $\mathbb{C}$ المعادلة ذات المجهول $z$ التالية: $z^2-6\sqrt2\,z+36=0$.

\textbf{2)} المستوي المركّب منسوب إلى المعلم المتعامد والمتجانس $(O;\vec u,\vec v)$. لتكن النقط $A$، $B$، $C$ و$D$ لواحقها على الترتيب:
\[z_A=3\sqrt2(1+i),\qquad z_B=\overline{z_A},\qquad z_C=6\sqrt2,\qquad z_D=\frac{z_C}{2}.\]
\begin{sq}
\qa اكتب $z_A$، $z_B$ و$(1+i)z_A$ على الشكل الأسّي.
\qb احسب $\left(\dfrac{(1+i)z_A}{6\sqrt2}\right)^{2014}$.
\qc بيّن أنّ النقط $O$، $A$، $B$ و$C$ تنتمي إلى نفس الدائرة التي مركزها $D$، يُطلب تعيين نصف قطرها.
\qd احسب $\dfrac{z_B-z_C}{z_A-z_C}$ ثم جد قيسًا للزاوية $(\overrightarrow{CA};\overrightarrow{CB})$. ما هي طبيعة الرباعي $OACB$؟
\end{sq}
\textbf{3)} ليكن $R$ الدوران الذي مركزه $O$ وزاويته $\dfrac{\pi}{2}$.
\begin{sq}
\qa اكتب العبارة المركّبة للدوران $R$.
\qb عيّن لاحقة النقطة $C'$ صورة $C$ بالدوران $R$، ثم تحقّق أنّ النقط $A$، $C$ و$C'$ في استقامية.
\qc عيّن لاحقة النقطة $A'$ صورة $A$ بالدوران $R$، ثم حدّد صورة الرباعي $OACB$ بالدوران $R$.
\end{sq}
\end{tamrin}

% ====================== التمرين: 2014 م2 ======================
\begin{tamrin}{بكالوريا 2014 -- الموضوع الثاني -- علوم تجريبية}
\textbf{1)} حلّ في مجموعة الأعداد المركّبة $\mathbb{C}$ المعادلة ذات المجهول $z$ التالية: $(z-i)(z^2-2z+5)=0$.

\textbf{2)} في المستوي المنسوب إلى المعلم المتعامد والمتجانس $(O;\vec u,\vec v)$ (وحدة الطول $1\,\mathrm{cm}$) تُعطى النقط $A$، $B$ و$C$ لواحقها على الترتيب: $z_A=i$، $z_B=1+2i$، $z_C=1-2i$.
\begin{sq}
\qa أنشئ النقط $A$، $B$ و$C$.
\qb جد $z_H$ لاحقة النقطة $H$ المسقط العمودي للنقطة $A$ على $(BC)$.
\qc احسب مساحة المثلث $ABC$.
\end{sq}
\textbf{3)} ليكن $S$ التشابه المباشر الذي مركزه $A$ ونسبته $\dfrac12$ وزاويته $\dfrac{\pi}{2}$.
\begin{sq}
\qa عيّن الكتابة المركّبة للتشابه $S$.
\qb بيّن أنّ مساحة صورة المثلث $ABC$ بالتشابه $S$ تساوي $\dfrac12\,\mathrm{cm}^2$.
\end{sq}
\textbf{4)} $M$ نقطة لاحقتها $z$. عيّن مجموعة النقط $M$ حيث: $|z|=|iz+1+2i|$.
\end{tamrin}

% ====================== التمرين: 2013 م1 ======================
\begin{tamrin}{بكالوريا 2013 -- الموضوع الأول -- علوم تجريبية}
\textbf{1)} حلّ في مجموعة الأعداد المركّبة $\mathbb{C}$ المعادلة $(I)$ ذات المجهول $z$ التالية:
\[z^2-(4\cos\alpha)z+4=0\qquad\ldots(I)\]
حيث $\alpha$ وسيط حقيقي.

\textbf{2)} من أجل $\alpha=\dfrac{\pi}{3}$، نرمز إلى حلّي المعادلة $(I)$ بـ $z_1$ و$z_2$. بيّن أنّ $\left(\dfrac{z_1}{z_2}\right)^{2013}=1$.

\textbf{3)} نعتبر في المستوي المنسوب إلى المعلم المتعامد والمتجانس $(O;\vec u,\vec v)$ النقط $A$، $B$ و$C$ التي لواحقها على الترتيب:
\[z_A=1+i\sqrt3,\qquad z_B=1-i\sqrt3,\qquad z_C=4+i\sqrt3 .\]
\begin{sq}
\qa أنشئ النقط $A$، $B$ و$C$.
\qb اكتب على الشكل الجبري العدد المركّب $\dfrac{z_C-z_A}{z_B-z_A}$، ثم استنتج أنّ $C$ هي صورة $B$ بالتشابه المباشر $S$ الذي مركزه $A$ ويُطلب تعيين نسبته وزاويته.
\qc عيّن لاحقة النقطة $G$ مرجح الجملة $\{(A;1),(B;-1),(C;2)\}$، ثم أنشئ $G$.
\qd احسب $z_D$ لاحقة النقطة $D$ بحيث يكون الرباعي $ABDG$ متوازي أضلاع.
\end{sq}
\end{tamrin}

% ====================== التمرين: 2013 م2 ======================
\begin{tamrin}{بكالوريا 2013 -- الموضوع الثاني -- علوم تجريبية}
نعتبر في مجموعة الأعداد المركّبة $\mathbb{C}$ المعادلة $(E)$ ذات المجهول $z$ الآتية:
\[z^2+4z+13=0\qquad\ldots(E)\]
\textbf{1)} تحقّق أنّ العدد المركّب $-2-3i$ حلّ للمعادلة $(E)$، ثم جد الحلّ الآخر.

\textbf{2)} $A$ و$B$ نقطتان من المستوي المركّب لاحقتاهما $z_A=-2-3i$ و$z_B=i$ على الترتيب.
$S$ التشابه المباشر الذي مركزه $A$، نسبته $\dfrac12$ وزاويته $\dfrac{\pi}{2}$ والذي يحوّل كلّ نقطة $M(z)$ من المستوي إلى النقطة $M'(z')$.
\begin{sq}
\qa بيّن أنّ: $z'=\dfrac12 iz-\dfrac72-2i$.
\qb احسب $z_C$ لاحقة النقطة $C$، علمًا أنّ $C$ هي صورة $B$ بالتشابه $S$.
\end{sq}
\textbf{3)} لتكن النقطة $D$ حيث: $2\overrightarrow{AD}+\overrightarrow{AB}=\vec 0$.
\begin{sq}
\qa بيّن أنّ $D$ مرجح النقطتين $A$ و$B$ المرفقتين بمعاملين حقيقيين يُطلب تعيينهما.
\qb احسب $z_D$ لاحقة النقطة $D$.
\qc بيّن أنّ $\dfrac{z_D-z_A}{z_C-z_A}=i$، ثم استنتج طبيعة المثلث $ACD$.
\end{sq}
\end{tamrin}

% ====================== التمرين: 2012 م1 ======================
\begin{tamrin}{بكالوريا 2012 -- الموضوع الأول -- علوم تجريبية}
\textbf{1)} نعتبر في مجموعة الأعداد المركّبة $\mathbb{C}$ المعادلة ذات المجهول $z$:
$z=\dfrac{3i(z+2i)}{z-2+3i}$ (حيث $z\neq2-3i$).
\begin{itemize}
\item حلّ في $\mathbb{C}$ هذه المعادلة.
\end{itemize}
\textbf{2)} ينسب المستوي المركّب إلى المعلم المتعامد والمتجانس $(O;\vec u,\vec v)$. $A$، $B$ نقطتان لاحقتاهما على الترتيب $z_A$ و$z_B$ حيث: $z_A=1+i\sqrt5$ و$z_B=1-i\sqrt5$.
\begin{itemize}
\item تحقّق أنّ $A$ و$B$ تنتميان إلى دائرة مركزها $O$ يُطلب تعيين نصف قطرها.
\end{itemize}
\textbf{3)} نرفق بكلّ نقطة $M$ من المستوي لاحقتها $z$ ($z\neq2-3i$) النقطة $M'$ لاحقتها $z'$ حيث: $z'=\dfrac{3i(z+2i)}{z-2+3i}$.
النقط $C$، $D$، $E$ لواحقها على الترتيب: $z_C=-2i$، $z_D=2-3i$ و$z_E=3i$، و$(\Delta)$ محور القطعة $[CD]$.
\begin{sq}
\qa عبّر عن المسافة $OM'$ بدلالة المسافتين $CM$ و$DM$.
\qb استنتج أنّه من أجل كلّ نقطة $M$ من $(\Delta)$ فإنّ النقطة $M'$ تنتمي إلى دائرة $(\gamma)$ يُطلب تعيين مركزها ونصف قطرها. تحقّق أنّ $E$ تنتمي إلى $(\gamma)$.
\end{sq}
\end{tamrin}

% ====================== التمرين: 2012 م2 ======================
\begin{tamrin}{بكالوريا 2012 -- الموضوع الثاني -- علوم تجريبية}
\textbf{1)} $P(z)$ كثير حدود للمتغيّر المركّب $z$ حيث: $P(z)=z^3-12z^2+48z-72$.
\begin{sq}
\qa تحقّق أنّ $6$ هو جذر لكثير الحدود $P(z)$.
\qb حدّد العددين الحقيقيين $\alpha$ و$\beta$ بحيث من أجل كلّ عدد مركّب $z$: $P(z)=(z-6)(z^2+\alpha z+\beta)$.
\qc حلّ في مجموعة الأعداد المركّبة $\mathbb{C}$ المعادلة $P(z)=0$.
\end{sq}
\textbf{2)} المستوي المركّب منسوب إلى المعلم المتعامد والمتجانس $(O;\vec u,\vec v)$. $A$، $B$ و$C$ نقط من المستوي المركّب لواحقها على الترتيب:
\[z_A=6,\qquad z_B=3+i\sqrt3,\qquad z_C=3-i\sqrt3 .\]
\begin{sq}
\qa اكتب كلًّا من $z_A$، $z_B$ و$z_C$ على الشكل الأسّي.
\qb اكتب العدد المركّب $\dfrac{z_A-z_B}{z_A-z_C}$ على الشكل الجبري، ثم على الشكل الأسّي.
\qc استنتج طبيعة المثلث $ABC$.
\end{sq}
\textbf{3)} ليكن $S$ التشابه المباشر الذي مركزه $C$، نسبته $\sqrt3$ وزاويته $\dfrac{\pi}{2}$.
\begin{sq}
\qa جد الكتابة المركّبة للتشابه $S$.
\qb عيّن $z_{A'}$ لاحقة النقطة $A'$ صورة النقطة $A$ بالتشابه $S$.
\qc بيّن أنّ النقط $A$، $B$ و$A'$ في استقامية.
\end{sq}
\end{tamrin}

% ====================== التمرين: 2011 م1 ======================
\begin{tamrin}{بكالوريا 2011 -- الموضوع الأول -- علوم تجريبية}
نعتبر في المستوي المنسوب إلى المعلم المتعامد والمتجانس $(O;\vec u,\vec v)$ النقط $A$، $B$ و$C$ لواحقها على الترتيب:
\[z_A=-i,\qquad z_B=2+3i,\qquad z_C=-4+i .\]
\textbf{1)}
\begin{sq}
\qa اكتب على الشكل الجبري العدد المركّب $\dfrac{z_C-z_A}{z_B-z_A}$.
\qb عيّن طويلة العدد المركّب $\dfrac{z_C-z_A}{z_B-z_A}$ وعمدة له، ثم استنتج طبيعة المثلث $ABC$.
\end{sq}
\textbf{2)} نعتبر التحويل النقطي $T$ في المستوي الذي يرفق بكلّ نقطة $M$ لاحقتها $z$ النقطة $M'$ ذات اللاحقة $z'$ بحيث: $z'=iz-1-i$.
\begin{sq}
\qa عيّن طبيعة التحويل $T$ محدّدًا عناصره المميّزة.
\qb ما هي صورة النقطة $B$ بالتحويل $T$؟
\end{sq}
\textbf{3)} لتكن $D$ ذات اللاحقة $z_D=-6+2i$.
\begin{sq}
\qa بيّن أنّ النقط $A$، $C$ و$D$ في استقامية.
\qb عيّن نسبة التحاكي $h$ الذي مركزه $A$ ويحوّل النقطة $C$ إلى $D$.
\qc عيّن العناصر المميّزة للتشابه المباشر $S$ الذي مركزه $A$ ويحوّل $B$ إلى $D$.
\end{sq}
\end{tamrin}

% ====================== التمرين: 2011 م2 ======================
\begin{tamrin}{بكالوريا 2011 -- الموضوع الثاني -- علوم تجريبية}
نعتبر في المستوي المنسوب إلى المعلم المتعامد والمتجانس $(O;\vec u,\vec v)$ النقط $A$، $B$ و$C$ لواحقها على الترتيب:
\[z_A=3-2i,\qquad z_B=3+2i,\qquad z_C=4i .\]
\textbf{1)}
\begin{sq}
\qa علّم النقط $A$، $B$ و$C$.
\qb ما هي طبيعة الرباعي $OABC$؟ علّل إجابتك.
\qc عيّن لاحقة النقطة $\Omega$ مركز الرباعي $OABC$.
\end{sq}
\textbf{2)} عيّن ثم أنشئ $(E)$ مجموعة النقط $M$ من المستوي التي تحقّق: $\left\|\overrightarrow{MO}+\overrightarrow{MA}+\overrightarrow{MB}+\overrightarrow{MC}\right\|=12$.

\textbf{3)}
\begin{sq}
\qa حلّ في مجموعة الأعداد المركّبة $\mathbb{C}$ المعادلة: $z^2-6z+13=0$. نسمّي $z_0$، $z_1$ حلّي هذه المعادلة.
\qb لتكن $M$ نقطة من المستوي لاحقتها العدد المركّب $z$. عيّن مجموعة النقط $M$ من المستوي التي تحقّق: $|z-z_0|=|z-z_1|$.
\end{sq}
\end{tamrin}

% ====================== التمرين: 2010 م1 ======================
\begin{tamrin}{بكالوريا 2010 -- الموضوع الأول -- علوم تجريبية}
نعتبر في المستوي المنسوب إلى المعلم المتعامد والمتجانس $(O;\vec u,\vec v)$ النقطتين $A$ و$B$ لاحقتاهما على الترتيب: $z_A=1+i$ و$z_B=3i$.

\textbf{1)} اكتب كلًّا من $z_A$ و$z_B$ على الشكل الأسّي.

\textbf{2)} ليكن $S$ التشابه المباشر الذي يرفق بكلّ نقطة $M$ لاحقتها $z$ النقطة $M'$ لاحقتها $z'$ بحيث: $z'=2iz+6+3i$.
\begin{sq}
\qa عيّن العناصر المميّزة للتشابه المباشر $S$.
\qb عيّن $z_C$ لاحقة النقطة $C$ صورة النقطة $A$ بالتشابه المباشر $S$.
\qc استنتج طبيعة المثلث $ABC$.
\end{sq}
\textbf{3)} لتكن النقطة $D$ مرجح الجملة $\{(A;2),(B;-2),(C;2)\}$.
\begin{sq}
\qa عيّن $z_D$ لاحقة النقطة $D$.
\qb عيّن مع التبرير طبيعة الرباعي $ABCD$.
\end{sq}
\textbf{4)} لتكن $M$ نقطة من المستوي تختلف عن $B$ وعن $D$ لاحقتها $z$، ولتكن $(\Delta)$ مجموعة النقط $M$ ذات اللاحقة $z$ التي يكون من أجلها $\dfrac{z_B-z}{z_D-z}$ عددًا حقيقيًا موجبًا تمامًا.
\begin{sq}
\qa تحقّق أنّ النقطة $E$ ذات اللاحقة $z_E=6+3i$ تنتمي إلى $(\Delta)$.
\qb أعط تفسيرًا هندسيًا لعمدة العدد المركّب $\dfrac{z_B-z}{z_D-z}$. عيّن حينئذ المجموعة $(\Delta)$.
\end{sq}
\end{tamrin}

% ====================== التمرين: 2010 م2 ======================
\begin{tamrin}{بكالوريا 2010 -- الموضوع الثاني -- علوم تجريبية}
\textbf{1)} حلّ في مجموعة الأعداد المركّبة $\mathbb{C}$ المعادلة: $z^2-6z+18=0$، ثم اكتب الحلّين على الشكل الأسّي.

\textbf{2)} في المستوي المنسوب إلى المعلم المتعامد والمتجانس $(O;\vec u,\vec v)$ نعتبر النقط $A$، $B$، $C$ و$D$ لواحقها على الترتيب:
\[z_A=3+3i,\qquad z_B=\overline{z_A},\qquad z_C=-z_A,\qquad z_D=-z_B .\]
\begin{sq}
\qa بيّن أنّ النقط $A$، $B$، $C$ و$D$ تنتمي إلى نفس الدائرة ذات المركز $O$.
\qb عيّن زاوية للدوران $R$ الذي مركزه $O$ ويحوّل النقطة $A$ إلى $B$.
\qc بيّن أنّ النقط $A$، $O$ و$C$ في استقامية وكذلك النقط $B$، $O$ و$D$.
\qd استنتج طبيعة الرباعي $ABCD$.
\end{sq}
\end{tamrin}

% ====================== التمرين: 2009 م1 ======================
\begin{tamrin}{بكالوريا 2009 -- الموضوع الأول -- علوم تجريبية}
$P(z)$ كثير حدود حيث: $P(z)=(z-1-i)(z^2-2z+4)$ و$z$ عدد مركّب.

\textbf{1)} حلّ في المجموعة $\mathbb{C}$ المعادلة: $P(z)=0$.

\textbf{2)} نضع: $z_1=1+i$، $z_2=1-i\sqrt3$.
\begin{sq}
\qa اكتب $z_1$ و$z_2$ على الشكل الأسّي.
\qb اكتب $\dfrac{z_1}{z_2}$ على الشكل الجبري ثم الشكل الأسّي.
\qc استنتج القيمة المضبوطة لكلّ من $\sin\left(\dfrac{7\pi}{12}\right)$ و$\cos\left(\dfrac{7\pi}{12}\right)$.
\end{sq}
\textbf{3)}
\begin{sq}
\qa $n$ عدد طبيعي، عيّن قيم $n$ بحيث يكون العدد $\left(\dfrac{z_1}{z_2}\right)^{n}$ حقيقيًا.
\qb احسب قيمة العدد $\left(\dfrac{z_1}{z_2}\right)^{456}$.
\end{sq}
\end{tamrin}

% ====================== التمرين: 2009 م2 ======================
\begin{tamrin}{بكالوريا 2009 -- الموضوع الثاني -- علوم تجريبية}
المستوي منسوب إلى المعلم المتعامد والمتجانس $(O;\vec i,\vec j)$.

\textbf{1)} حلّ في المجموعة $\mathbb{C}$ المعادلة: $z^2-2z+4=0$.

\textbf{2)} نسمّي $z_1$ و$z_2$ حلّي هذه المعادلة.
\begin{sq}
\qa اكتب العددين $z_1$ و$z_2$ على الشكل الأسّي.
\qb $A$، $B$ و$C$ هي النقط من المستوي التي لواحقها على الترتيب: $z_A=1-i\sqrt3$، $z_B=1+i\sqrt3$ و$z_C=\dfrac12\left(5+i\sqrt3\right)$.
احسب الأطوال $AB$، $AC$ و$BC$، ثم استنتج طبيعة المثلث $ABC$.
\qc جد الطويلة وعمدة للعدد المركّب $z$ حيث: $z=\dfrac{z_C-z_B}{z_A-z_B}$.
\qd احسب $z^3$، $z^6$ ثم استنتج أنّ $z^{3k}$ عدد حقيقي من أجل كلّ عدد طبيعي $k$.
\end{sq}
\end{tamrin}

% ====================== التمرين: 2008 م1 ======================
\begin{tamrin}{بكالوريا 2008 -- الموضوع الأول -- علوم تجريبية}
\textbf{1)} حلّ في مجموعة الأعداد المركّبة $\mathbb{C}$ المعادلة: $z^2-(1+2i)z-1+i=0$.
نرمز للحلّين بـ $z_1$ و$z_2$ حيث $|z_1|<|z_2|$. بيّن أنّ $\left(\dfrac{z_1}{z_2}\right)^{2008}$ عدد حقيقي.

\textbf{2)} المستوي المركّب منسوب إلى المعلم المتعامد والمتجانس $(O;\vec u,\vec v)$.
لتكن $A$، $B$ و$C$ نقطًا من المستوي لواحقها على الترتيب: $1$، $z_1$ و$z_2$. ليكن $z$ العدد المركّب حيث: $z=\dfrac{z_2-1}{z_1-1}$.
\begin{sq}
\qa انطلاقًا من التعريف $e^{i\theta}=\cos\theta+i\sin\theta$ ومن الخاصية $e^{i(\theta_1+\theta_2)}=e^{i\theta_1}\times e^{i\theta_2}$،
برهن أنّ: $e^{-i\theta}=\dfrac{1}{e^{i\theta}}$ وأنّ: $\dfrac{e^{i\theta_1}}{e^{i\theta_2}}=e^{i(\theta_1-\theta_2)}$ حيث $\theta$، $\theta_1$ و$\theta_2$ أعداد حقيقية.
\qb اكتب $z$ على الشكل الأسّي.
\qc اكتب $z$ على الشكل المثلثي واستنتج أنّ $C$ هي صورة $B$ بتشابه مباشر مركزه $A$، يُطلب تعيين نسبته وزاويته.
\end{sq}
\end{tamrin}

% ====================== التمرين: 2008 م2 ======================
\begin{tamrin}{بكالوريا 2008 -- الموضوع الثاني -- علوم تجريبية}
\textbf{1)} حلّ في مجموعة الأعداد المركّبة $\mathbb{C}$ المعادلة ذات المجهول $z$ التالية: $z^2+iz-2-6i=0$.

\textbf{2)} نعتبر في المستوي المركّب المنسوب إلى المعلم المتعامد والمتجانس $(O;\vec u,\vec v)$ النقطتين $A$ و$B$ اللتين لاحقتاهما $z_A$ و$z_B$ على الترتيب حيث: $z_A=2+i$ و$z_B=-2-2i$.
عيّن $z_\omega$ لاحقة النقطة $\omega$ مركز الدائرة $(\Gamma)$ ذات القطر $[AB]$.

\textbf{3)} لتكن $C$ النقطة ذات اللاحقة $z_C$ حيث: $z_C=\dfrac{4-i}{1+i}$.
اكتب $z_C$ على الشكل الجبري ثم أثبت أنّ النقطة $C$ تنتمي إلى الدائرة $(\Gamma)$.

\textbf{4)}
\begin{sq}
\qa برهن أنّ عبارة التشابه المباشر الذي مركزه $M_0(z_0)$ ونسبته $k$ ($k>0$) وزاويته $\theta$ والذي يرفق بكلّ نقطة $M(z)$ النقطة $M'(z')$ هي: $z'-z_0=ke^{i\theta}(z-z_0)$.
\qb \textbf{تطبيق:} عيّن الطبيعة والعناصر المميّزة للتحويل $S$ المعرّف بـ: $z'+\dfrac12 i=2e^{i\frac{\pi}{3}}\left(z+\dfrac12 i\right)$.
\end{sq}
\end{tamrin}

\section{تمارين البكالوريا -- شعبة تقني رياضي}
% ---------- helpers for the figures of this part (complex plane) ----------
% \begin{cpoCP}{u}{v}{xmin}{xmax}{ymin}{ymax}{unit cm} ... \end{cpoCP}
\newenvironment{cpoCP}[7]{\par\begin{center}\begin{LTR}\begin{tikzpicture}[x=#7cm,y=#7cm,font=\footnotesize,line cap=round]
 \draw[gray!20,very thin] (#3,#5) grid (#4,#6);
 \draw[-{Stealth[length=2mm]},gray!80!black] (#3,0)--(#4,0);
 \draw[-{Stealth[length=2mm]},gray!80!black] (0,#5)--(0,#6);
 \draw[very thick,-{Stealth[length=1.8mm]},black] (0,0)--(1,0) node[below,pos=0.55]{$\vec{#1}$};
 \draw[very thick,-{Stealth[length=1.8mm]},black] (0,0)--(0,1) node[left,pos=0.55]{$\vec{#2}$};
 \node[below left,inner sep=1.5pt] at (0,0){$O$};}
{\end{tikzpicture}\end{LTR}\end{center}}
% filled point with label: \cpoD{coord}{anchor}{label}
\newcommand{\cpoD}[3]{\fill (#1) circle (1.7pt) node[#2]{$#3$};}
% right-angle mark at #1 between directions #2 and #3
\newcommand{\cpoRA}[3]{\draw[thin] ($(#1)!2.6mm!(#2)$) -- ($($(#1)!2.6mm!(#2)$)+($(#1)!2.6mm!(#3)$)-(#1)$) -- ($(#1)!2.6mm!(#3)$);}
% equal-length tick on segment [#1#2]
\newcommand{\cpoT}[2]{\path ($(#1)!0.5!(#2)$) coordinate (cpoM); \draw[thick] ($(cpoM)!1.5mm!90:(#2)$) -- ($(cpoM)!1.5mm!-90:(#2)$);}
\newcommand{\cpoTT}[2]{\path ($(#1)!0.5!(#2)$) coordinate (cpoM); \draw[thick] ($($(cpoM)!1.5mm!90:(#2)$)!0.6mm!(#2)$) -- ($($(cpoM)!1.5mm!-90:(#2)$)!0.6mm!(#2)$); \draw[thick] ($($(cpoM)!1.5mm!90:(#2)$)!-0.6mm!(#2)$) -- ($($(cpoM)!1.5mm!-90:(#2)$)!-0.6mm!(#2)$);}

% ===================== 2019 - S1 - TM =====================
\begin{tamrin}{بكالوريا 2019 -- الموضوع الأول -- تقني رياضي}
المستوي المركب منسوب إلى المعلم المتعامد والمتجانس $(O;\vec u,\vec v)$. $A$، $B$ و $C$ النقط التي لواحقها على الترتيب:
\[z_A=1+i,\qquad z_B=2+i,\qquad z_C=\frac32+i\left(1+\frac{\sqrt3}{2}\right).\]
$(\Gamma)$ الدائرة التي مركزها $A$ ونصف قطرها $1$.

\QQ{1}
\begin{sq}
\qa تحقق أن النقطة $C$ من الدائرة $(\Gamma)$.
\qb عين قيسا بالراديان للزاوية $\left(\overrightarrow{AB};\overrightarrow{AC}\right)$، ثم استنتج أن $C$ صورة $B$ بالدوران $r$ الذي مركزه $A$ يطلب تعيين زاويته.
\end{sq}
\QQ{2} $S$ التشابه المباشر الذي يحول النقطة $M$ ذات اللاحقة $z$ إلى النقطة $M'$ ذات اللاحقة $z'$ حيث: $z'=(1+i\sqrt3)z+\sqrt3-i\sqrt3$.
\begin{sq}
\qa حدد العناصر المميزة للتشابه $S$.
\qb عين $z_D$ لاحقة $D$ صورة $B$ بالتشابه $S$.
\end{sq}
\QQ{3} ما هي نسبة التحاكي $h$ الذي مركزه $A$ حيث: $S=h\circ r$؟ استنتج أن النقط $A$، $C$ و $D$ في استقامية.

\QQ{4} $(E)$ مجموعة النقط $M$ من المستوي ذات اللاحقة $z$ حيث: $z=z_A+k\,e^{i\frac{\pi}{3}}$ مع $k\in\mathbb{R}_+^*$.
تحقق أن النقطة $C$ تنتمي إلى $(E)$، ثم حدد طبيعة المجموعة $(E)$.
\end{tamrin}

% ===================== 2019 - S2 - TM =====================
\begin{tamrin}{بكالوريا 2019 -- الموضوع الثاني -- تقني رياضي}
\textbf{I)} تحقق أن: $(2-2\sqrt3)^2=16-8\sqrt3$.
\textbf{II)} المستوي المركب منسوب إلى المعلم المتعامد والمتجانس $(O;\vec u,\vec v)$. نعتبر النقط $A$، $B$ و $C$ التي لواحقها على الترتيب:
\[z_A=4e^{i\frac{\pi}{3}}+4e^{i\frac{5\pi}{6}},\qquad z_B=\frac12iz_A,\qquad z_C=-\frac14z_A.\]
\QQ{1} أكتب $z_A$ على الشكل الجبري، ثم بين أن: $z_A=4\sqrt2\,e^{i\frac{7\pi}{12}}$.

\QQ{2} استنتج القيمتين المضبوطتين للعددين الحقيقيين $\cos\left(\dfrac{7\pi}{12}\right)$ و $\sin\left(\dfrac{7\pi}{12}\right)$.

\QQ{3} $S$ التشابه المباشر الذي يحول $A$ إلى $B$ و $B$ إلى $C$. لتكن $M'$ النقطة ذات اللاحقة $z'$ صورة النقطة $M$ ذات اللاحقة $z$ بالتشابه $S$.
\begin{sq}
\qa بين أن: $z'=\dfrac12iz$.
\qb حدد العناصر المميزة للتشابه $S$.
\end{sq}
\QQ{4} $G$ النقطة ذات اللاحقة $z_G$ مرجح الجملة $\{(A;2),(B;-2),(C;4)\}$.
\begin{sq}
\qa بين أن: $z_G=2e^{i\frac{\pi}{3}}$.
\qb $(E)$ مجموعة النقط $M$ من المستوي ذات اللاحقة $z$ حيث: $\left\|\overrightarrow{MA}-\overrightarrow{MB}+2\overrightarrow{MC}\right\|=2\sqrt2$.
حدد طبيعة المجموعة $(E)$ وعناصرها المميزة، ثم أحسب محيط $(E')$ صورة $(E)$ بالتشابه $S$.
\end{sq}
\end{tamrin}

% ===================== 2018 - S1 - TM =====================
\begin{tamrin}{بكالوريا 2018 -- الموضوع الأول -- تقني رياضي}
\textbf{I)} حل في مجموعة الأعداد المركبة $\mathbb{C}$ المعادلة ذات المجهول $z$: $z^2-2\sqrt2z+4=0$.

\textbf{II)} المستوي المركب منسوب إلى المعلم المتعامد والمتجانس $(O;\vec u,\vec v)$. لتكن النقطتان $A$ و $B$ لاحقتاهما: $z_A=\sqrt2+i\sqrt2$ و $z_B=\overline{z_A}$.

\QQ{1} أكتب على الشكل الأسي كلا من العددين المركبين $z_A$ و $\dfrac{1}{z_B}$، ثم بين أن العدد $\left(\dfrac{2}{z_B}\right)^{2018}$ تخيلي صرف.

\QQ{2} لتكن النقطة $C$ صورة النقطة $B$ بالتحاكي $h$ الذي مركزه $\omega$ ذات اللاحقة $z_\omega=\dfrac{\sqrt2}{2}$ ونسبته $-3$.
بين أن لاحقة النقطة $C$ هي: $z_C=-\sqrt2+3\sqrt2\,i$.

\QQ{3} عين لاحقة النقطة $D$ صورة النقطة $B$ بالدوران $R$ الذي مركزه $O$ وزاويته $-\dfrac{\pi}{2}$.

\QQ{4}
\begin{sq}
\qa بين أن $\dfrac{z_C-z_A}{z_D-z_A}=-i$، ثم استنتج طبيعة المثلث $ACD$.
\qb جد لاحقة النقطة $E$ بحيث يكون الرباعي $ACED$ مربعا.
\end{sq}
\end{tamrin}

% ===================== 2018 - S2 - TM =====================
\begin{tamrin}{بكالوريا 2018 -- الموضوع الثاني -- تقني رياضي}
\textbf{I)}
\begin{sq}
\qa حل في مجموعة الأعداد المركبة $\mathbb{C}$ المعادلة ذات المجهول $z$: $4z^2-2z+1=0\quad\dots(E)$.
\qb أكتب العددين $\dfrac{1}{z_1}$ و $\dfrac{1}{z_2}$ على الشكل الأسي، حيث $z_1$ و $z_2$ حلا المعادلة $(E)$.
\end{sq}
\textbf{II)} المستوي المركب منسوب إلى المعلم المتعامد والمتجانس $(O;\vec u,\vec v)$. نعتبر النقط $A$، $B$ و $C$ التي لواحقها على الترتيب: $z_A=4$، $z_B=1+i\sqrt3$ و $z_C=1-i\sqrt3$.

\QQ{1}
\begin{sq}
\qa أحسب $\dfrac{z_B-z_A}{z_C-z_A}$، ثم حدد طبيعة المثلث $ABC$.
\qb استنتج أن $B$ هي صورة $C$ بدوران مركزه $A$ يطلب تعيين زاويته.
\end{sq}
\QQ{2} جد لاحقة النقطة $D$ صورة النقطة $A$ بالانسحاب الذي شعاعه $\overrightarrow{CB}$، واستنتج بدقة طبيعة الرباعي $ACBD$.

\QQ{3} حدد طبيعة $(\gamma)$ مجموعة النقط $M$ من المستوي ذات اللاحقة $z$ التي تحقق: $\left|iz+\sqrt3-i\right|=\left|z-1+i\sqrt3\right|$.

\QQ{4} بين أن النقطة $G$ مركز الدائرة المحيطة بالمثلث $ABC$ تنتمي إلى $(\gamma)$.
\end{tamrin}

% ===================== 2017 R - S1 - TM =====================
\begin{tamrin}{بكالوريا 2017 (الدورة الاستدراكية) -- الموضوع الأول -- تقني رياضي}
\textbf{I)} حل في مجموعة الأعداد المركبة $\mathbb{C}$ المعادلة ذات المجهول $z$: $(z-4)(z^2-2z+4)=0$.

\textbf{II)} المستوي المركب منسوب إلى المعلم المتعامد والمتجانس $(O;\vec u,\vec v)$. نعتبر النقط $A$، $B$ و $C$ التي لواحقها على الترتيب: $z_A=4$، $z_B=1+i\sqrt3$ و $z_C=1-i\sqrt3$.

\QQ{1} أكتب العدد المركب $\dfrac{z_C-z_A}{z_B-z_A}$ على الشكل الأسي، ثم استنتج طبيعة المثلث $ABC$.

\QQ{2}
\begin{sq}
\qa عين لاحقة النقطة $D$ صورة النقطة $B$ بالدوران $R$ الذي مركزه $O$ وزاويته $\dfrac{2\pi}{3}$.
\qb عين طبيعة الرباعي $ABDC$.
\end{sq}
\QQ{3} من أجل كل عدد طبيعي $n$ نضع: $z_n=(z_B)^n+(z_C)^n$.
\begin{sq}
\qa بين أنه من أجل كل عدد طبيعي $n$: $z_n=2^{n+1}\times\cos\left(\dfrac{n\pi}{3}\right)$.
\qb نضع من أجل كل عدد طبيعي $n$: $t_n=z_{6n}$. عبر عن $t_n$ بدلالة $n$، ثم أحسب $P_n$ بدلالة $n$ حيث: $P_n=t_0\times t_1\times t_2\times\dots\times t_n$.
\end{sq}
\end{tamrin}

% ===================== 2017 R - S2 - TM =====================
\begin{tamrin}{بكالوريا 2017 (الدورة الاستدراكية) -- الموضوع الثاني -- تقني رياضي}
\textbf{I)} حل في مجموعة الأعداد المركبة $\mathbb{C}$ المعادلة ذات المجهول $z$: $(z+1-\sqrt3)(z^2+2z+4)=0$.

\textbf{II)} المستوي المركب منسوب إلى المعلم المتعامد والمتجانس $(O;\vec u,\vec v)$. نعتبر النقط $A$، $B$ و $C$ التي لواحقها على الترتيب: $z_A=-1+\sqrt3$، $z_B=-1-i\sqrt3$ و $z_C=\overline{z_B}$.

\QQ{1} بين أن: $z_B-z_A=i(z_C-z_A)$، ثم استنتج طبيعة المثلث $ABC$ وأحسب مساحته.

\QQ{2}
\begin{sq}
\qa أكتب على الشكل الجبري العدد $L$ حيث: $L=\dfrac{z_C-z_A}{z_C}$.
\qb بين أن: $L=\dfrac{\sqrt6}{2}\left(\cos\dfrac{\pi}{12}+i\sin\dfrac{\pi}{12}\right)$، ثم استنتج القيمة المضبوطة لـ $\tan\dfrac{\pi}{12}$.
\end{sq}
\QQ{3} نعتبر التحويل النقطي $S$ الذي يحول النقطة $M$ ذات اللاحقة $z$ إلى النقطة $M'$ ذات اللاحقة $z'$ والمعرف بـ: $z'=(z-z_B)L+z_B$.
بين أن $S$ تشابه مباشر يطلب تحديد عناصره المميزة.

\QQ{4} لتكن النقط $A'$، $B'$ و $C'$ صور النقط $A$، $B$ و $C$ على الترتيب بالتحويل $S\circ S$. أحسب مساحة المثلث $A'B'C'$.
\end{tamrin}

% ===================== 2017 - S1 - TM =====================
\begin{tamrin}{بكالوريا 2017 -- الموضوع الأول -- تقني رياضي}
المستوي المركب منسوب إلى المعلم المتعامد والمتجانس $(O;\vec u,\vec v)$. نعتبر النقط $A$، $B$ و $C$ التي لواحقها على الترتيب: $z_A=-1$، $z_B=2+i$ و $z_C=-i$.

\QQ{1} أكتب العدد المركب $\dfrac{z_A-z_C}{z_B-z_C}$ على الشكل الأسي، ثم استنتج طبيعة المثلث $ABC$.

\QQ{2} عين العبارة المركبة للتشابه المباشر $S$ الذي مركزه $C$ ويحول $B$ إلى $A$.

\QQ{3} نعتبر النقطة $D$ نظيرة $B$ بالنسبة إلى $C$ والنقطة $E$ صورة النقطة $D$ بالتشابه $S$.
\begin{sq}
\qa عين $z_D$ لاحقة النقطة $D$، ثم تحقق أن: $z_E=1-2i$ حيث $z_E$ لاحقة النقطة $E$.
\qb حدد طبيعة الرباعي $ADEB$.
\end{sq}
\QQ{4} $(\Gamma)$ مجموعة النقط $M$ من المستوي ذات اللاحقة $z$ ($M$ تختلف عن $A$ و $B$) حيث:
$\arg(z-z_A)-\arg(z-z_B)=\dfrac{\pi}{2}+2k\pi$ مع $k\in\mathbb{Z}$.
تحقق أن النقطة $C$ تنتمي إلى $(\Gamma)$، ثم حدد طبيعة المجموعة $(\Gamma)$ وأنشئها.
\end{tamrin}

% ===================== 2017 - S2 - TM =====================
\begin{tamrin}{بكالوريا 2017 -- الموضوع الثاني -- تقني رياضي}
المستوي المركب منسوب إلى المعلم المتعامد والمتجانس $(O;\vec u,\vec v)$. نعتبر النقط $A$، $B$، $C$ و $D$ التي لواحقها على الترتيب:
\[z_A=1+i,\qquad z_B=\overline{z_A},\qquad z_C=\frac12(1-i),\qquad z_D=\overline{z_C}.\]
\QQ{1}
\begin{sq}
\qa أكتب $z_A$ و $z_C$ على الشكل الأسي، ثم استنتج الشكل الأسي لكل من $z_B$ و $z_D$.
\qb عين قيم العدد الطبيعي $n$ التي تحقق: $(z_A)^n=(z_B)^n$.
\end{sq}
\QQ{2}
\begin{sq}
\qa جد نسبة ومركز التحاكي $h$ الذي يحول $A$ إلى $D$ ويحول $B$ إلى $C$.
\qb أحسب طويلة العدد المركب $\dfrac{z_C-z_B}{z_D-z_A}$ ثم استنتج طبيعة الرباعي $ADCB$.
\end{sq}
\QQ{3} أحسب $z_G$ لاحقة النقطة $G$ مرجح الجملة $\{(A,2);(B,2);(C,-1);(D,-1)\}$.

\QQ{4} لتكن $(\Gamma)$ مجموعة النقط $M$ من المستوي بحيث: $\left\|2\overrightarrow{MA}+2\overrightarrow{MB}-\overrightarrow{MC}-\overrightarrow{MD}\right\|=\sqrt5$.
بين أن النقطة $A$ تنتمي إلى $(\Gamma)$، ثم حدد طبيعة المجموعة $(\Gamma)$ وعناصرها المميزة وأنشئها.
\end{tamrin}

% ===================== 2016 - S1 - TM =====================
\begin{tamrin}{بكالوريا 2016 -- الموضوع الأول -- تقني رياضي}
\QQ{1} حل في المجموعة $\mathbb{C}$ المعادلة ذات المجهول $z$: $9z^2-6\sqrt3z+4=0$.

\QQ{2} في المستوي المركب المنسوب إلى المعلم المتعامد والمتجانس $(O;\vec u,\vec v)$، لتكن النقطتان $A$ و $B$ لاحقتاهما على الترتيب: $z_A=\dfrac{\sqrt3}{3}+\dfrac13i$ و $z_B=\overline{z_A}$.
\begin{sq}
\qa أكتب كلا من $z_A$ و $z_B$ على الشكل الأسي.
\qb بين أن: $\left(\dfrac{z_A}{z_B}\right)^{2016}+\left(\dfrac{z_A}{z_B}\right)^{1437}=0$.
\qc عين قيم العدد الطبيعي $n$ بحيث يكون $\left(\dfrac{z_A}{z_B}\right)^n$ عددا حقيقيا.
\end{sq}
\QQ{3} $f$ التحويل النقطي الذي يرفق بكل نقطة $M$ لاحقتها $z$ النقطة $M'$ لاحقتها $z'$ حيث: $z'=\left(\dfrac{z_A}{z_B}\right)z$.
\begin{sq}
\qa عين طبيعة التحويل النقطي $f$ وعناصره المميزة.
\qb أحسب $z_C$ لاحقة النقطة $C$ صورة النقطة $A$ بالتحويل $f$.
\qc عين $z_D$ لاحقة النقطة $D$ حتى تكون $O$ مركز ثقل الرباعي $ABCD$.
\end{sq}
\end{tamrin}

% ===================== 2016 - S2 - TM =====================
\begin{tamrin}{بكالوريا 2016 -- الموضوع الثاني -- تقني رياضي}
\QQ{1}
\begin{sq}
\qa حل في المجموعة $\mathbb{C}$ المعادلة ذات المجهول $z$: $(2z-\sqrt2)(z^2-2\sqrt2z+4)=0$.
\qb أكتب الحلول على الشكل الأسي.
\end{sq}
\QQ{2} المستوي المركب منسوب إلى المعلم المتعامد والمتجانس $(O;\vec u,\vec v)$. $A$، $B$ و $C$ النقط التي لواحقها على الترتيب: $a=\dfrac{\sqrt2}{2}$، $b=\sqrt2+i\sqrt2$ و $c=\sqrt2-i\sqrt2$.
\begin{sq}
\qa علم النقط $A$، $B$ و $C$ في المعلم السابق.
\qb نعتبر النقطة $D$ صورة النقطة $C$ بالتشابه المباشر $S$ الذي مركزه $A$ ونسبته $3$ وزاويته $\pi$، والنقطة $E$ صورة النقطة $C$ بالدوران $R$ الذي مركزه $O$ وزاويته $-\dfrac{\pi}{2}$.
أحسب اللاحقتين $d$ و $e$ للنقطتين $D$ و $E$ على الترتيب.
\end{sq}
\QQ{3} نضع: $z=\dfrac{d-b}{e-b}$.
\begin{sq}
\qa أكتب العدد $z$ على الشكل المثلثي.
\qb نعتبر النقطة $I$ منتصف القطعة $[DE]$، و $F$ نظيرة النقطة $B$ بالنسبة للنقطة $I$. ما طبيعة الرباعي $BDFE$؟
\end{sq}
\end{tamrin}

% ===================== 2015 - S1 - TM =====================
\begin{tamrin}{بكالوريا 2015 -- الموضوع الأول -- تقني رياضي}
نعتبر في المستوي المنسوب إلى المعلم المتعامد والمتجانس $(O;\vec u,\vec v)$ النقطتين $A$ و $B$ اللتين لاحقتاهما على الترتيب $z_A$ و $z_B$ حيث: $z_A=1-i$ و $z_B=3+3i$.

\QQ{1}
\begin{sq}
\qa أكتب $z_A$ و $z_B$ على الشكل الأسي.
\qb $n$ عدد طبيعي، عين قيم $n$ بحيث يكون العدد $\left(\dfrac{z_A}{\sqrt2}\right)^n$ حقيقيا.
\qc $z$ عدد مركب حيث $\dfrac{z}{z_A}=4e^{i\frac{\pi}{12}}$، أحسب طويلة العدد $z$ وعمدة له، ثم أكتب $\dfrac{z}{z_A}$ على الشكل الجبري.
\qd استنتج $\cos\dfrac{\pi}{12}$ و $\sin\dfrac{\pi}{12}$.
\end{sq}
\QQ{2}
\begin{sq}
\qa أحسب $z_C$ لاحقة النقطة $C$ صورة النقطة $B$ بالدوران الذي مركزه $A$ وزاويته $\dfrac{\pi}{2}$، ثم استنتج طبيعة المثلث $ABC$.
\qb أحسب $z_D$ لاحقة النقطة $D$ مرجح الجملة $\{(A,-1);(B,1);(C,1)\}$، ثم بين أن $ABDC$ مربع.
\end{sq}
\end{tamrin}

% ===================== 2015 - S2 - TM =====================
\begin{tamrin}{بكالوريا 2015 -- الموضوع الثاني -- تقني رياضي}
\QQ{1} حل في مجموعة الأعداد المركبة $\mathbb{C}$ المعادلة ذات المجهول $z$: $z^2-4(\sin\theta)z+4=0\quad\dots(I)$، حيث $\theta$ وسيط حقيقي.

\QQ{2} من أجل $\theta=\dfrac{\pi}{3}$ نرمز إلى حلي المعادلة $(I)$ بـ $z_1$ و $z_2$. أكتب $z_1$ و $z_2$ على الشكل الأسي.

\QQ{3} نعتبر في المستوي المركب المنسوب إلى المعلم المتعامد والمتجانس $(O;\vec u,\vec v)$ النقط $A$، $B$ و $C$ التي لواحقها على الترتيب: $z_A=\sqrt3+i$، $z_B=\sqrt3-i$ و $z_C=3\sqrt3+i$.
\begin{sq}
\qa أكتب العدد المركب $\dfrac{z_C-z_A}{z_B-z_A}$ على الشكل الجبري، ثم على الشكل الأسي. استنتج طبيعة المثلث $ABC$.
\qb استنتج أن النقطة $C$ هي صورة النقطة $B$ بالتشابه المباشر $S$ الذي مركزه $A$ يطلب تعيين نسبته وزاوية له.
\qc عين $z_D$ لاحقة النقطة $D$ صورة النقطة $B$ بالانسحاب الذي شعاعه $\overrightarrow{AC}$، ثم حدد طبيعة الرباعي $ABDC$.
\end{sq}
\QQ{4}
\begin{sq}
\qa عين $(\Gamma_1)$ مجموعة النقط $M$ من المستوي ذات اللاحقة $z$ حيث: $\dfrac{z-z_C}{z-z_B}$ تخيلي صرف مع $z\neq z_B$.
\qb عين $(\Gamma_2)$ مجموعة النقط $M$ من المستوي ذات اللاحقة $z$ حيث: $\dfrac{z-z_C}{z-z_B}$ حقيقي مع $z\neq z_B$.
\end{sq}
\end{tamrin}

% ===================== 2014 - S1 - TM =====================
\begin{tamrin}{بكالوريا 2014 -- الموضوع الأول -- تقني رياضي}
\QQ{1} حل في مجموعة الأعداد المركبة $\mathbb{C}$ المعادلة ذات المجهول $z$: $(z-i)(z^2-2\sqrt3z+4)=0$.

\QQ{2} المستوي المركب منسوب إلى المعلم المتعامد والمتجانس $(O;\vec u,\vec v)$. $A$، $B$ و $C$ النقط التي لواحقها على الترتيب: $z_1=\sqrt3+i$، $z_2=\sqrt3-i$ و $z_3=i$.
\begin{sq}
\qa أكتب العدد $\dfrac{z_1}{z_2}$ على الشكل الأسي.
\qb هل توجد قيم للعدد الطبيعي $n$ يكون من أجلها العدد المركب $\left(\dfrac{z_1}{z_2}\right)^n$ تخيليا صرفا؟ برر إجابتك.
\end{sq}
\QQ{3}
\begin{sq}
\qa أعط العبارة المركبة للتشابه $S$ الذي مركزه $A$ ويحول $B$ إلى $C$، محددا نسبته وزاويته.
\qb استنتج طبيعة المثلث $ABC$.
\end{sq}
\QQ{4}
\begin{sq}
\qa عين العناصر المميزة لـ $(E)$ مجموعة النقط $M$ من المستوي ذات اللاحقة $z$ والتي تحقق: $|z-z_1|^2+|z-z_3|^2=5$.
\qb عين $(E')$ مجموعة النقط $M$ من المستوي التي لاحقتها $z$ حيث: $|z-z_1|=|z-z_3|$.
\end{sq}
\end{tamrin}

% ===================== 2014 - S2 - TM =====================
\begin{tamrin}{بكالوريا 2014 -- الموضوع الثاني -- تقني رياضي}
نعتبر في المستوي المركب المنسوب إلى المعلم المتعامد والمتجانس $(O;\vec u,\vec v)$ النقطة $A$ ذات اللاحقة $z_0=1+i$.

\QQ{1}
\begin{sq}
\qa عين ثم أنشئ $(\gamma)$ مجموعة النقط $M(z)$ من المستوي حيث: $z=z_0+2e^{i\theta}$ و $\theta$ يمسح $\mathbb{R}$.
\qb عين ثم أنشئ $(\gamma')$ مجموعة النقط $M(z)$ من المستوي حيث: $z=z_0+k\,e^{i\frac{3\pi}{4}}$ و $k$ يمسح $\mathbb{R}^+$.
\qc عين إحداثيات نقطة تقاطع $(\gamma)$ و $(\gamma')$.
\end{sq}
\QQ{2} نسمي $B$ النقطة التي لاحقتها $z_1$ حيث: $z_1=z_0+2e^{i\frac{3\pi}{4}}$.
\begin{sq}
\qa عين الشكل الجبري للعدد المركب $\dfrac{z_1-z_0}{z_0}$، ثم استنتج طبيعة المثلث $OAB$.
\qb عين $z_2$ لاحقة النقطة $C$ صورة النقطة $B$ بالدوران الذي مركزه $A$ وزاويته $-\dfrac{\pi}{2}$.
\qc عين العددين $\alpha$ و $\beta$ بحيث تكون النقطة $O$ مرجحا للجملة $\{(A,\alpha);(C,\beta)\}$ و $\alpha+\beta=\sqrt2$.
\qd عين ثم أنشئ $(E)$ مجموعة النقط $M$ من المستوي حيث: $\left((1+\sqrt2)\overrightarrow{MA}-\overrightarrow{MC}\right)\cdot\left(\overrightarrow{MA}-\overrightarrow{MC}\right)=0$.
\end{sq}
\end{tamrin}

% ================= 2013 S1 TM =================
\begin{tamrin}{بكالوريا 2013 -- الموضوع الأول -- تقني رياضي}
\QQ{1} حل في مجموعة الأعداد المركبة $\mathbb{C}$ المعادلة ذات المجهول $z$: $2z^2+6z+17=0$.
\QQ{2} المستوي المركب منسوب إلى المعلم المتعامد والمتجانس $(O;\vec u,\vec v)$.
$A$، $B$ و $C$ النقط التي لواحقها على الترتيب: $z_A=-4$، $z_B=-\dfrac32+\dfrac52 i$ و $z_C=-\dfrac32-\dfrac52 i$.

أحسب الطويلة وعمدة للعدد المركب $\dfrac{z_B-z_A}{z_C-z_A}$، ثم استنتج طبيعة المثلث $ABC$.
\QQ{3}
\begin{sq}
\qa عيّن $z_D$ و $z_E$ لاحقتي النقطتين $D$ و $E$ على الترتيب حتى يكون الرباعي $BCDE$ مربعًا مركزه $A$.
\qb عيّن $(\Gamma_1)$ مجموعة النقط $M$ من المستوي حيث: $\left\|\overrightarrow{MB}+\overrightarrow{MC}+\overrightarrow{MD}+\overrightarrow{ME}\right\|=10\sqrt2$.
\end{sq}
\QQ{4} $(\Gamma_2)$ مجموعة النقط $M$ من المستوي، ذات اللاحقة $z$ حيث: $\arg(z+4)=\dfrac{\pi}{4}$.
تحقق أن النقطة $B$ تنتمي إلى $(\Gamma_2)$، ثم عيّن المجموعة $(\Gamma_2)$.
\end{tamrin}

% ================= 2013 S2 TM =================
\begin{tamrin}{بكالوريا 2013 -- الموضوع الثاني -- تقني رياضي}
\QQ{1} حل في مجموعة الأعداد المركبة $\mathbb{C}$ المعادلة ذات المجهول $z$: $(z+5-i\sqrt3)(z^2+2z+4)=0$.
\QQ{2} المستوي المركب منسوب إلى المعلم المتعامد والمتجانس $(O;\vec u,\vec v)$.
$A$، $B$ و $C$ النقط التي لواحقها على الترتيب: $z_A=-1-i\sqrt3$، $z_B=-1+i\sqrt3$ و $z_C=-5+i\sqrt3$.
$S$ التشابه المباشر الذي يحوّل $A$ إلى $C$ ويحوّل $O$ إلى $B$.

جد العبارة المركبة للتشابه $S$، ثم عيّن العناصر المميزة له.
\QQ{3}
\begin{sq}
\qa عيّن $z_D$ لاحقة النقطة $D$ مرجح الجملة $\{(A,2);(B,-1);(C,1)\}$.
\qb أكتب العدد المركب $\dfrac{z_B-z_A}{z_D-z_A}$ على الشكل الأسي، ثم استنتج طبيعة المثلث $ABD$.
\qc عيّن المجموعة $(\Gamma)$ للنقط $M$ من المستوي حيث: $\left\|2\overrightarrow{MA}-\overrightarrow{MB}+\overrightarrow{MC}\right\|=\left\|\overrightarrow{MA}-\overrightarrow{MB}\right\|$.
\end{sq}
\end{tamrin}

% ================= 2012 S1 TM =================
\begin{tamrin}{بكالوريا 2012 -- الموضوع الأول -- تقني رياضي}
\QQ{1} عيّن العددين المركبين $z_1$ و $z_2$ بحيث: $\begin{cases}2z_1+3z_2=9-2i\\ 3z_1-z_2=8+8i\end{cases}$.
\QQ{2} نعتبر في المستوي المنسوب إلى المعلم المتعامد والمتجانس $(O;\vec u,\vec v)$ النقط $A$، $B$ و $\Omega$ التي لواحقها على الترتيب:
$z_A=3+2i$، $z_B=-3$ و $z_\Omega=1-2i$.
\begin{sq}
\qa أثبت أن: $z_B-z_\Omega=i(z_A-z_\Omega)$.
\qb عيّن طبيعة المثلث $\Omega AB$.
\end{sq}
\QQ{3} $h$ التحاكي الذي مركزه النقطة $A$ ونسبته $2$.
\begin{sq}
\qa عيّن الكتابة المركبة للتحاكي $h$.
\qb عيّن $z_C$ لاحقة النقطة $C$ صورة النقطة $\Omega$ بالتحاكي $h$.
\qc عيّن $z_D$ لاحقة النقطة $D$ مرجح الجملة $\{(A,1);(B,-1);(C,1)\}$.
\qd بيّن أن $ABCD$ مربع.
\end{sq}
\QQ{4} $(E)$ مجموعة النقط $M$ من المستوي التي تحقق: $\left\|\overrightarrow{MA}-\overrightarrow{MB}+\overrightarrow{MC}\right\|=4\sqrt5$.
\begin{sq}
\qa تحقق أن النقطة $B$ تنتمي إلى المجموعة $(E)$، ثم عيّن طبيعة $(E)$ وعناصرها المميزة.
\qb أنشئ المجموعة $(E)$.
\end{sq}
\end{tamrin}

% ================= 2012 S2 TM =================
\begin{tamrin}{بكالوريا 2012 -- الموضوع الثاني -- تقني رياضي}
\QQ{1} حل في مجموعة الأعداد المركبة $\mathbb{C}$ المعادلة ذات المجهول $z$: $(z^2+2z+4)(z^2-2\sqrt3z+4)=0$.
\QQ{2} المستوي المركب منسوب إلى المعلم المتعامد والمتجانس $(O;\vec u,\vec v)$.
$A$، $B$، $C$ و $D$ النقط التي لواحقها على الترتيب: $z_A=\sqrt3+i$، $z_B=\sqrt3-i$، $z_C=-1-\sqrt3 i$ و $z_D=-1+\sqrt3 i$.
\begin{sq}
\qa أكتب كلًا من $z_A$، $z_B$، $z_C$ و $z_D$ على الشكل الأسي.
\qb تحقق أن: $\dfrac{z_D-z_B}{z_A-z_C}=i$، ثم استنتج أن المستقيمين $(AC)$ و $(BD)$ متعامدان.
\end{sq}
\QQ{3} $z_n$ العدد المركب الذي طويلته $\dfrac{1}{2^n}$ و $\dfrac{2\pi}{3}n$ عمدة له حيث $n$ عدد طبيعي.
$L_n$ العدد المركب المعرّف بـ: $L_n=z_D\times z_n$.
\begin{sq}
\qa أكتب كلًا من $L_0$ و $L_1$ على الشكل الجبري.
\qb $(u_n)$ المتتالية المعرفة من أجل كل عدد طبيعي $n$ بـ: $u_n=|L_n|$.
\begin{itemize}
\item أثبت أن المتتالية $(u_n)$ هندسية يطلب تعيين أساسها وحدها الأول.
\item $M_0$، $M_1$، $\ldots$، $M_n$ صور الأعداد المركبة $L_0$، $L_1$، $\ldots$، $L_n$ على الترتيب.
أحسب، بدلالة $n$، المجموع $S_n$ حيث: $S_n=\left\|\overrightarrow{OM_1}\right\|+\left\|\overrightarrow{OM_2}\right\|+\cdots+\left\|\overrightarrow{OM_n}\right\|$.
\item جد نهاية $S_n$ لما يؤول $n$ إلى $+\infty$.
\end{itemize}
\end{sq}
\end{tamrin}

% ================= 2011 S1 TM =================
\begin{tamrin}{بكالوريا 2011 -- الموضوع الأول -- تقني رياضي}
نعتبر في مجموعة الأعداد المركبة $\mathbb{C}$ المعادلة: \quad $(E)\;\ldots\;z^2-2\sqrt3z+4=0$.
\QQ{1} حل في $\mathbb{C}$ المعادلة $(E)$، ثم أكتب حلولها على الشكل المثلثي.
\QQ{2} المستوي منسوب إلى المعلم المتعامد والمتجانس $(O;\vec u,\vec v)$.
$A$، $B$ و $C$ النقط التي لواحقها $z_A=2i$، $z_B=\sqrt3+i$ و $z_C=\sqrt3-i$ على الترتيب. نضع: $L=-\dfrac12-\dfrac{\sqrt3}{2}i$.
\begin{sq}
\qa أكتب $L$ على الشكل الأسي.
\qb أثبت أن $z_A-z_B=L(z_C-z_B)$، ثم استنتج أن $A$ صورة $C$ بتحويل نقطي يطلب تعيينه وتحديد عناصره المميزة.
\qc استنتج نوع المثلث $ABC$ ثم أحسب مساحته.
\end{sq}
\end{tamrin}

% ================= 2011 S2 TM =================
\begin{tamrin}{بكالوريا 2011 -- الموضوع الثاني -- تقني رياضي}
المستوي منسوب إلى المعلم المتعامد والمتجانس $(O;\vec u,\vec v)$. $L$ العدد المركب المعرّف كما يلي: $L=\dfrac{-4\sqrt2+i\sqrt2}{5+3i}$.
\QQ{1}
\begin{sq}
\qa أكتب $L$ على الشكل الجبري ثم على الشكل الأسي.
\qb بيّن أن $L^{12}+1=0$، ثم أحسب: $(-4\sqrt2+i\sqrt2)^{12}+(5+3i)^{12}$.
\qc $n$ عدد طبيعي فردي و $p$ عدد طبيعي زوجي. أثبت أن: $L^{4n}+L^{4p}=0$.
\end{sq}
\QQ{2}
\begin{sq}
\qa النقطتان $A$ و $B$ لاحقتاهما على الترتيب: $z_A=5+3i$ و $z_B=5-3i$.
عيّن اللاحقة $z_{A'}$ للنقطة $A'$ صورة النقطة $A$ بالتشابه المباشر الذي مركزه النقطة $B$ ونسبته $\sqrt2$ وزاويته $\dfrac{3\pi}{4}$.
\qb عيّن $z_G$ لاحقة النقطة $G$ مركز ثقل المثلث $ABA'$.
\end{sq}
\end{tamrin}

% ================= 2010 S1 TM =================
\begin{tamrin}{بكالوريا 2010 -- الموضوع الأول -- تقني رياضي}
\QQ{1} حل في مجموعة الأعداد المركبة $\mathbb{C}$ المعادلة ذات المجهول $z$: $(z-3+2i)(z^2+6z+10)=0$.
\QQ{2} علّم في المستوي المنسوب إلى المعلم المتعامد والمتجانس $(O;\vec u,\vec v)$ النقط $A$، $C$، $D$ و $I$ لواحقها:
$z_A=3-2i$، $z_C=-3+i$، $z_D=-3-i$ و $z_I=1$ على الترتيب.
\QQ{3} $z$ عدد مركب يحقق الجملة: $\begin{cases}\arg(z-3+2i)=\arg(z-1)+\dfrac{\pi}{2}\\[4pt] |z-3+2i|=|z-1|\end{cases}$.
\begin{sq}
\qa بيّن أن الجملة تكافئ: $\dfrac{z-3+2i}{z-1}=i$، ثم عيّن قيمة $z$.
\qb $B$ النقطة التي لاحقتها $z_B=3$، تحقق أن $\overrightarrow{AB}=\overrightarrow{DC}$. ما هي طبيعة الرباعي $ABCD$؟
\qc لتكن $J$ النقطة التي لاحقتها $z_J$ حيث: $z_J=1-2i$.
أكتب على الشكل الأسي العدد المركب $Z$ حيث: $Z=\dfrac{z_A-z_I}{z_B-z_J}$.
تحقق أن $\overrightarrow{AB}=\overrightarrow{JI}$. ما هي طبيعة الرباعي $ABIJ$؟
\end{sq}
\end{tamrin}

% ================= 2010 S2 TM =================
\begin{tamrin}{بكالوريا 2010 -- الموضوع الثاني -- تقني رياضي}
\QQ{1}
\begin{sq}
\qa أكتب على الشكل الأسي العدد المركب $a$ حيث: $a=-2+2i\sqrt3$.
\qb حل في مجموعة الأعداد المركبة $\mathbb{C}$ المعادلة ذات المجهول $z$: $z^2=-2+2i\sqrt3$.
\end{sq}
\QQ{2} ينسب المستوي إلى المعلم المتعامد والمتجانس $(O;\vec u,\vec v)$.
$A$، $B$ و $C$ النقط التي لواحقها $z_A=-2$، $z_B=-1-\sqrt3 i$ و $z_C=1+\sqrt3 i$ على الترتيب.
\begin{sq}
\qa أحسب طويلة العدد المركب $\dfrac{z_C-z_A}{z_B-z_A}$ وعمدة له.
\qb استنتج طبيعة المثلث $ABC$.
\end{sq}
\QQ{3} لتكن $(E)$ مجموعة النقط ذات اللاحقة $z$ حيث: $\arg(\overline{z}+2)=\dfrac{\pi}{3}$.
\begin{sq}
\qa تحقق أن $B$ تنتمي إلى $(E)$.
\qb عيّن المجموعة $(E)$.
\end{sq}
\end{tamrin}

% ================= 2009 S1 TM =================
\begin{tamrin}{بكالوريا 2009 -- الموضوع الأول -- تقني رياضي}
\QQ{1}
\begin{sq}
\qa حل في $\mathbb{C}$ المعادلة ذات المجهول $z$: $z^2-2z+2=0$.
\qb استنتج في $\mathbb{C}$ حلول المعادلة ذات المجهول $z$: $(\overline{z}+3)^2-2(\overline{z}+3)+2=0$، حيث $\overline{z}$ مرافق $z$.
\end{sq}
\QQ{2} المستوي المركب منسوب إلى المعلم المتعامد والمتجانس $(O;\vec u,\vec v)$.
النقط $A$، $B$ و $M$ لواحقها $(1+i)$، $(1-i)$ و $z$ على الترتيب.
\begin{sq}
\qa عيّن $(\Gamma)$ مجموعة النقط $M$ من المستوي حيث: $z=1-i+k\,e^{i\frac{5\pi}{4}}$ عندما $k$ يمسح $\R^+$.
\qb عيّن $(E)$ مجموعة النقط $M$ من المستوي حيث: $|z-1+i|=|z-1-i|$.
\end{sq}
\end{tamrin}

% ================= 2009 S2 TM =================
\begin{tamrin}{بكالوريا 2009 -- الموضوع الثاني -- تقني رياضي}
\QQ{1} حل في مجموعة الأعداد المركبة $\mathbb{C}$ المعادلة ذات المجهول $z$: $z^2-6z+18=0$.
\QQ{2} ليكن العدد المركب $z_1$ حيث: $z_1=3-3i$.
\begin{sq}
\qa أكتب $z_1$ على الشكل الأسي.
\qb أحسب طويلة العدد $z_3$ وعمدة له حيث: $z_1\times z_3=6\left(\cos\dfrac{\pi}{12}+i\sin\dfrac{\pi}{12}\right)$. استنتج قيمتي $\cos\dfrac{\pi}{12}$ و $\sin\dfrac{\pi}{12}$.
\end{sq}
\QQ{3} نعتبر في المستوي المركب المنسوب إلى المعلم المتعامد والمتجانس $(O;\vec u,\vec v)$ النقط $A$، $B$ و $C$ لواحقها $3+3i$، $3-3i$ و $\dfrac{\sqrt2}{2}+i\dfrac{\sqrt6}{2}$ على الترتيب.
\begin{sq}
\qa عيّن قيم العدد الحقيقي $\alpha$ حتى تقبل الجملة المثقلة $\{(A,1);(B,-1);(C,\alpha)\}$ مرجحًا نرمز له بـ $G_\alpha$.
\qb عيّن مجموعة النقط $G_\alpha$ لما يتغير $\alpha$ في $\R^*$.
\end{sq}
\end{tamrin}

% ================= 2008 S1 TM =================
\begin{tamrin}{بكالوريا 2008 -- الموضوع الأول -- تقني رياضي}
لتكن في مجموعة الأعداد المركبة $\mathbb{C}$ المعادلة $(*)$ المعرفة كما يلي:
\[(*)\;\ldots\quad z^3+(2-4i)z^2-(6+9i)z+9(-1+i)=0\]
\QQ{1} بيّن أن $z_0=3i$ هو حل للمعادلة $(*)$.
\QQ{2} لتكن $A$، $B$ و $C$ صور الحلول $z_0$، $z_1$ و $z_2$ على الترتيب في المستوي المنسوب إلى المعلم المتعامد والمتجانس $(O;\vec u,\vec v)$.
عيّن النقطة $G$ مرجح الجملة $\{(A,1);(B,1);(C,-1)\}$.
\QQ{3} عيّن المجموعة $(E)$ للنقط $M$ حيث: $AM^2+BM^2-CM^2=-13$.
بيّن أن النقطة $A$ تنتمي إلى المجموعة $(E)$، ثم أنشئ $(E)$.
\QQ{4} تحقق أن النقط $O$، $B$ و $G$ في استقامية، ثم عيّن صورة المجموعة $(E)$ بالتحاكي الذي مركزه النقطة $O$ ويحوّل $B$ إلى $G$ محددًا عناصره المميزة.
\end{tamrin}

% ================= 2008 S2 TM =================
\begin{tamrin}{بكالوريا 2008 -- الموضوع الثاني -- تقني رياضي}
$r$ عدد حقيقي موجب تمامًا و $\theta$ عدد حقيقي كيفي.
\QQ{1} حل في مجموعة الأعداد المركبة $\mathbb{C}$ المعادلة ذات المجهول $z$:
\[z^2-2i\left(r\cos\frac{\theta}{2}\right)z-r^2=0\]
أكتب الحلين على الشكل الأسي.
\QQ{2} في المستوي المركب المنسوب إلى المعلم المتعامد والمتجانس $(O;\vec u,\vec v)$ نعتبر النقطتين $A$ و $B$ صورتي الحلين.
عيّن $\theta$ حتى يكون المثلث $OAB$ متقايس الأضلاع.
\end{tamrin}

\section{تمارين البكالوريا -- شعبة رياضيات}
% ===================== التمرين 1 =====================
\begin{tamrin}{بكالوريا 2026 -- الموضوع الأول -- شعبة رياضيات}
\QQ{1} حل في المجموعة $\mathbb{C}$ المعادلة ذات المجهول $z$ التالية: $\left(z-1+i\sqrt3\right)\left(z^2-2z+2\right)=0$.

\QQ{2} في المستوي المركب المنسوب إلى المعلم المتعامد والمتجانس $\left(O;\vec u,\vec v\right)$، نعتبر النقط $A$، $B$، $C$ و$D$ لاحقاتها على الترتيب $z_A$، $z_B$، $z_C$ و$z_D$ حيث:
\[z_A=1+i,\qquad z_B=1-i\sqrt3,\qquad z_C=1-i,\qquad z_D=\frac{1-\sqrt3}{2}(1+i).\]
\begin{sq}
\qa أكتب كلا من $z_A$، $z_B$، $z_C$ و$z_D$ على الشكل المثلثي.
\qb استنتج الشكل المثلثي للعدد المركب $L$ حيث: $L=\dfrac{z_A}{z_B}$.
\qc أكتب العدد $L$ على الشكل الجبري ثم استنتج القيمة المضبوطة لكل من $\cos\dfrac{7\pi}{12}$ و$\sin\dfrac{7\pi}{12}$.
\end{sq}
\QQ{3}
\begin{sq}
\qa تحقق أن: $z_D-z_A=i\left(z_D-z_B\right)$ ثم بين أن المثلث $ABD$ قائم ومتساوي الساقين.
\qb عين $z_E$ لاحقة النقطة $E$ مرجح الجملة $\left\{(A;1),(B;1),(D;-1)\right\}$ ثم حدد طبيعة الرباعي $AEBD$.
\end{sq}
\end{tamrin}

% ===================== التمرين 2 =====================
\begin{tamrin}{بكالوريا 2026 -- الموضوع الثاني -- شعبة رياضيات}
$\theta$ عدد حقيقي حيث: $-\dfrac{\pi}{2}<\theta<\dfrac{\pi}{2}$.

\QQ{1} حل في المجموعة $\mathbb{C}$ المعادلة ذات المجهول $z$ التالية: $\left(z-i\sqrt3\right)\left(z^2-4(\sin\theta)z+4\right)=0$.

\QQ{2} في المستوي المركب المنسوب إلى المعلم المتعامد والمتجانس $\left(O;\vec u,\vec v\right)$، نعتبر النقط $A$، $B$، $C$ و$D$ لاحقاتها على الترتيب $z_A$، $z_B$، $z_C$ و$z_D$ حيث:
\[z_A=i\sqrt3,\quad z_B=2(\sin\theta+i\cos\theta),\quad z_C=2(\sin\theta-i\cos\theta),\quad z_D=1+i\sqrt3.\]
\begin{sq}
\qa أكتب كلا من $z_A$، $z_B$، $z_C$ و$z_D$ على الشكل المثلثي.
\qb استنتج أن النقط $B$، $C$ و$D$ تنتمي إلى نفس الدائرة التي يطلب تعيين مركزها ونصف قطرها.
\qc بين أن: $\left(\dfrac{z_A}{\sqrt3}\right)^{2026}+\left(\dfrac{z_A}{\sqrt3}\right)^{1447}=-1-i$.
\qd عين قيم العدد $\theta$ حتى يكون المثلث $ABC$ قائما في $B$.
\end{sq}
\QQ{3} نضع $\theta=\dfrac{\pi}{6}$ ونعتبر العدد المركب $L$ حيث: $L=-(1+i)z_B$.
\begin{sq}
\qa أكتب $L$ على الشكل المثلثي وعلى الشكل الجبري.
\qb استنتج القيمة المضبوطة لكل من $\cos\dfrac{19\pi}{12}$ و$\sin\dfrac{19\pi}{12}$.
\end{sq}
\end{tamrin}

% ===================== التمرين 3 =====================
\begin{tamrin}{بكالوريا 2025 -- الموضوع الأول -- شعبة رياضيات}
\QQ{1} حل في مجموعة الأعداد المركبة $\mathbb{C}$ المعادلة ذات المجهول $z$ التالية: $\left(z^2+1\right)\left(z^2-6z+13\right)=0$.

\QQ{2} المستوي المركب منسوب إلى المعلم المتعامد والمتجانس $\left(O;\vec u,\vec v\right)$.
$A$، $B$ و$C$ نقط من المستوي لاحقاتها على الترتيب $z_A$، $z_B$ و$z_C$ حيث: $z_A=2-3i$، $z_B=iz_A$ و$z_C=-i$،
و$L$ عدد مركب حيث: $L=\dfrac{z_A-z_C}{z_B-z_C}$.
\begin{sq}
\qa أكتب كلا من العددين $z_A-z_C$ و$z_B-z_C$ على الشكل المثلثي.
\qb عين الطويلة وعمدة للعدد المركب $L$ ثم حدد طبيعة المثلث $ABC$.
\qc بين أن النقط $O$، $A$، $B$ و$C$ تنتمي إلى نفس الدائرة التي يطلب تعيين مركزها ونصف قطرها.
\qd من أجل كل عدد طبيعي $n$، نضع: $S_n=|L|+|L|^2+\cdots+|L|^n$.

أحسب بدلالة $n$ المجموع $S_n$ ثم استنتج $\displaystyle\lim_{n\to+\infty}S_n$.
\end{sq}
\QQ{3} $z_\theta$ عدد مركب حيث: $z_\theta=1+\cos2\theta-i\sin2\theta$ و$\theta$ عدد حقيقي يختلف عن $\dfrac{\pi}{2}$ مع $0<\theta<\pi$.

بين أن: $z_\theta=2\cos\theta\left(\cos\theta-i\sin\theta\right)$ ثم استنتج حسب قيم $\theta$ الطويلة وعمدة للعدد $z_\theta$.
\end{tamrin}

% ===================== التمرين 4 =====================
\begin{tamrin}{بكالوريا 2025 -- الموضوع الثاني -- شعبة رياضيات}
\QQ{1} حل في مجموعة الأعداد المركبة $\mathbb{C}$ المعادلة ذات المجهول $z$ التالية: $\left(iz-1+2i\right)\left(z^2-4z+5\right)=0$.

\QQ{2} المستوي المركب منسوب إلى المعلم المتعامد والمتجانس $\left(O;\vec u,\vec v\right)$.
$A$، $B$ و$C$ نقط من المستوي لاحقاتها على الترتيب $z_A$، $z_B$ و$z_C$ حيث: $z_A=2-i$، $z_B=\overline{z_A}$ و$z_C=-2-i$،
و$L$ عدد مركب حيث: $L=\dfrac{z_C-z_A}{z_B-z_A}$.
\begin{sq}
\qa أكتب العدد $L$ على الشكل المثلثي ثم حدد طبيعة المثلث $ABC$.
\qb بين أن النقط $A$، $B$ و$C$ تنتمي إلى نفس الدائرة التي يطلب تعيين مركزها ونصف قطرها.
\qc عين $z_D$ و$z_E$ لاحقتي النقطتين $D$ و$E$ على الترتيب حتى يكون الرباعي $BCDE$ معينا مركزه $A$.
\end{sq}
\QQ{3} $z_\theta$ عدد مركب حيث: $z_\theta=1-\cos2\theta+i\sin2\theta$ و$\theta$ عدد حقيقي غير معدوم مع $-\dfrac{\pi}{2}<\theta<\dfrac{\pi}{2}$.

بين أن: $z_\theta=2\sin\theta\left(\sin\theta+i\cos\theta\right)$ ثم استنتج حسب قيم $\theta$ الطويلة وعمدة للعدد $z_\theta$.
\end{tamrin}

% ===================== التمرين 5 =====================
\begin{tamrin}{بكالوريا 2024 -- الموضوع الأول -- شعبة رياضيات}
\textbf{I)} حل في مجموعة الأعداد المركبة $\mathbb{C}$ المعادلة ذات المجهول $z$ الآتية: $\left(z-8+6i\right)\left(z^2-2z+4\right)=0$.
\textbf{II)} المستوي المركب منسوب إلى المعلم المتعامد والمتجانس $\left(O;\vec i,\vec j\right)$.
نعتبر النقط $A$، $B$ و$C$ التي لاحقاتها على الترتيب $z_A$، $z_B$ و$z_C$ حيث: $z_A=1+i\sqrt3$، $z_B=iz_A$ و$z_C=-z_A$.

\QQ{1} تحقق أن: $z_A-z_B=i\left(z_C-z_B\right)$ ثم بين أن المثلث $ABC$ قائم ومتساوي الساقين.

\QQ{2}
\begin{sq}
\qa أكتب كلا من $z_A$، $z_B$ و$z_C$ على الشكل المثلثي.
\qb استنتج أن النقط $A$، $B$ و$C$ تنتمي إلى نفس الدائرة، يطلب تعيين مركزها ونصف قطرها.
\end{sq}
\QQ{3} النقطة $D$ هي نظيرة $B$ بالنسبة إلى مبدأ المعلم. بين أن الرباعي $ABCD$ مربع.
\end{tamrin}

% ===================== التمرين 6 =====================
\begin{tamrin}{بكالوريا 2024 -- الموضوع الثاني -- شعبة رياضيات}
\textbf{I)} حل في مجموعة الأعداد المركبة $\mathbb{C}$ المعادلة ذات المجهول $z$ الآتية: $\left(z^2+2i\right)\left(z^2-2\sqrt3z+6\right)=0$.

\textbf{II)} المستوي المركب منسوب إلى المعلم المتعامد والمتجانس $\left(O;\vec i,\vec j\right)$.
نعتبر النقط $A$، $B$ و$C$ التي لاحقاتها على الترتيب $z_A$، $z_B$ و$z_C$ حيث: $z_A=1-i$، $z_B=-z_A$ و$z_C=\sqrt3(1+i)$.

\QQ{1} أكتب كلا من $z_A$، $z_B$ و$z_C$ على الشكل المثلثي.

\QQ{2} أكتب العدد المركب $\dfrac{z_C-z_A}{z_B-z_A}$ على الشكل الجبري ثم المثلثي وبين أن المثلث $ABC$ متقايس الأضلاع.

\QQ{3}
\begin{sq}
\qa عين لاحقة النقطة $G$ مركز الدائرة المحيطة بالمثلث $ABC$ ثم أحسب نصف قطرها.
\qb النقطة $D$ هي نظيرة $C$ بالنسبة إلى مبدأ المعلم. بين أن الرباعي $ACBD$ معين.
\end{sq}
\end{tamrin}
\begin{molahada}
في نص الموضوع: «بين أن الرباعي $ABCD$ معين». لكن $O$ منتصف $[AB]$ ومنتصف $[CD]$، فقطرا المعين هما $[AB]$ و$[CD]$، والترتيب $A$، $B$، $C$، $D$ يعطي رباعيا متقاطع الأضلاع؛ لذلك صححنا التسمية إلى $ACBD$.
\end{molahada}

% ===================== التمرين 7 =====================
\begin{tamrin}{بكالوريا 2023 -- الموضوع الثاني -- شعبة رياضيات}
\QQ{1} حل في مجموعة الأعداد المركبة $\mathbb{C}$ المعادلة ذات المجهول $z$: $\left(\overline{z}-1+i\sqrt3\right)\left(z^2-2\sqrt2z+4\right)=0$.

\QQ{2} المستوي المركب منسوب إلى المعلم المتعامد والمتجانس $\left(O;\vec u,\vec v\right)$.
نعتبر النقط $A$، $B$ و$C$ التي لاحقاتها على الترتيب $z_A$، $z_B$ و$z_C$ حيث: $z_A=\sqrt2(1+i)$، $z_B=\overline{z_A}$ و$z_C=1+i\sqrt3$.
\begin{sq}
\qa أكتب $z_A$، $z_B$ و$z_C$ على الشكل المثلثي.
\qb استنتج أن النقط $A$، $B$ و$C$ تنتمي إلى نفس الدائرة التي يطلب تعيين مركزها ونصف قطرها.
\end{sq}
\QQ{3} نضع: $K=\dfrac{z_C}{2z_A}$.
\begin{sq}
\qa أحسب طويلة العدد المركب $K$ وعمدة له ثم أكتبه على الشكل الجبري.
\qb استنتج القيمة المضبوطة لكل من $\cos\dfrac{\pi}{12}$ و$\sin\dfrac{\pi}{12}$.
\end{sq}
\QQ{4} $n$ عدد طبيعي، نضع: $L_n=z_A^{\,n}+z_B^{\,n}$. بين أنه من أجل كل عدد طبيعي $n$، العدد المركب $L_n$ حقيقي.
\end{tamrin}

% ===================== التمرين 8 =====================
\begin{tamrin}{بكالوريا 2019 -- الموضوع الأول -- شعبة رياضيات}
في المستوي المركب المنسوب إلى المعلم المتعامد والمتجانس $\left(O;\vec u,\vec v\right)$، نعتبر النقط $A$، $B$، $C$، $D$ و$E$ حيث:
\[z_A=1+i\sqrt2,\quad z_B=i,\quad z_C=\overline{z_A},\quad z_D=\overline{z_B},\quad z_E=1.\]
\QQ{1} حل في المجموعة $\mathbb{C}$ المعادلة ذات المجهول $z$: $\left(z^2+1\right)\left(z^2-2z+3\right)=0$.

\QQ{2}
\begin{sq}
\qa أحسب كلا من $\left|z_A-1\right|$، $\left|z_B-1\right|$ و$\left|z_C-z_E\right|$ ثم تحقق أن النقط الأربعة $A$، $B$، $C$ و$D$ تنتمي إلى نفس الدائرة التي يطلب تعيين مركزها وطول نصف قطرها.
\qb بين أن: $z_B-z_E=\dfrac{\sqrt2}{2}(1+i)\left(z_A-z_E\right)$، ثم استنتج أن $B$ هي صورة $A$ بتحويل نقطي يطلب تعيين عناصره المميزة. ما طبيعة المثلث $ABE$؟
\end{sq}
\QQ{3} عين لاحقتي الشعاعين $\overrightarrow{BD}$ و$\overrightarrow{AE}$ محددا طبيعة الرباعي $ABDE$.

\QQ{4} $\vec w_1$ و$\vec w_2$ شعاعان من المستوي لاحقتاهما على الترتيب $z_1$ و$z_2$.
\begin{sq}
\qa برهن أن: ($\vec w_1$ و$\vec w_2$ متعامدان) يكافئ ($z_1\overline{z_2}+\overline{z_1}z_2=0$).
\qb عين مجموعة النقط $M$ من المستوي ذات اللاحقة $z$ حيث: $\left(z-z_A\right)\left(\overline{z}-z_D\right)+\left(z-z_B\right)\left(\overline{z}-z_C\right)=0$.
\end{sq}
\end{tamrin}

% ===================== التمرين 9 =====================
\begin{tamrin}{بكالوريا 2019 -- الموضوع الثاني -- شعبة رياضيات}
\QQ{1} نضع من أجل كل عدد مركب $z$: $P(z)=z^4-6z^3+29z^2-24z+100$.
\begin{sq}
\qa بين أنه من أجل كل عدد مركب $z$: $\overline{P(z)}=P\left(\overline{z}\right)$، ثم استنتج أنه إذا كان $z$ حلا للمعادلة $P(z)=0$ فإن $\overline{z}$ حل لها.
\qb حل في مجموعة الأعداد المركبة $\mathbb{C}$ المعادلة $P(z)=0$ علما أنها تقبل حلا تخيليا صرفا.
\end{sq}
\QQ{2} نعتبر في المستوي المركب المنسوب إلى المعلم المتعامد والمتجانس $\left(O;\vec u,\vec v\right)$ النقط $A$، $B$، $M$ و$M'$ التي لاحقاتها على الترتيب: $2i$، $3-4i$، $z$ و$z'$ حيث: $z'=\dfrac{-iz+4+3i}{z-2i}$ مع $z\neq2i$.
ولتكن $I$ مرجح الجملة $\left\{(A;2),(B;1)\right\}$ و$J$ مرجح الجملة $\left\{(A;-2),(B;1)\right\}$.
\begin{sq}
\qa عين اللاحقتين $z_I$ و$z_J$ للنقطتين $I$ و$J$ على الترتيب.
\qb لتكن $(E)$ مجموعة النقط $M(z)$ التي يكون من أجلها $|z'|=2$.

بين أن ($M$ من $(E)$) يكافئ ($\overrightarrow{IM}\cdot\overrightarrow{JM}=0$)، ثم عين $(E)$ وأنشئها.
\qc لتكن $(\Gamma)$ مجموعة النقط $M(z)$ التي يكون من أجلها $\arg\left(z'\right)=2k\pi$ حيث $k$ عدد صحيح.

تحقق أن النقطة $D$ ذات اللاحقة $\dfrac92-\dfrac52i$ تنتمي إلى $(\Gamma)$، ثم عين وأنشئ $(\Gamma)$.
\end{sq}
\QQ{3} عين الشكل الجبري للاحقة النقطة $G$ تقاطع المجموعتين $(E)$ و$(\Gamma)$.
\end{tamrin}

% ===================== التمرين 10 =====================
\begin{tamrin}{بكالوريا 2018 -- الموضوع الأول -- شعبة رياضيات}
المستوي منسوب إلى المعلم المتعامد والمتجانس $\left(O;\vec u,\vec v\right)$، $\theta$ عدد حقيقي من المجال $\left]-\pi;\pi\right]$.

\textbf{I)} حل في $\mathbb{C}$ المعادلة ذات المجهول $z$: $\left(z^2-2z+2\right)\left(z^2-2(\sin\theta)z+1\right)=0$.

\textbf{II)} $A$، $B$، $C$ و$D$ نقط من المستوي لاحقاتها على الترتيب:
$z_A=-\sqrt2\,e^{i\frac{5\pi}{4}}$، $z_B=1-i$، $z_C=\sin\theta+i\cos\theta$ و$z_D=\overline{z_C}$ ($\overline{z_C}$ يرمز إلى مرافق $z_C$).

\QQ{1} أكتب الأعداد $z_A$، $z_B$، $z_C$ و$z_D=\overline{z_C}$ على الشكل الأسي.

\QQ{2} $E$ نقطة من المستوي لاحقتها $z_E$ حيث: $z_E=\dfrac{z_A}{z_B}$. بين أن النقط $C$، $D$ و$E$ تنتمي إلى دائرة واحدة يطلب تعيين مركزها ونصف قطرها.

\QQ{3} ليكن $S$ التشابه المباشر الذي مركزه النقطة $A$ وزاويته $\dfrac{\pi}{4}$ ونسبته $\left(2\sqrt2-2\right)$.
عين قيمة $\theta$ حتى تكون النقطة $B$ صورة النقطة $C$ بالتشابه المباشر $S$.

\QQ{4} نضع: $\theta=-\dfrac{3\pi}{4}$. عين قيم العدد الطبيعي $n$ والتي من أجلها يكون العدد $\left(z_D\right)^n$ تخيليا صرفا.
\end{tamrin}

% ===================== التمرين 11 =====================
\begin{tamrin}{بكالوريا 2018 -- الموضوع الثاني -- شعبة رياضيات}
\QQ{1} $m$ عدد حقيقي. نعتبر في $\mathbb{C}$ المعادلة $(E)$ التالية: $z^2+(m+1)z+(2m-1)=0$.
عين قيم العدد الحقيقي $m$ التي من أجلها تقبل المعادلة $(E)$ حلين مركبين غير حقيقيين.

\QQ{2} نضع: $m=3$. حل المعادلة $(E)$.

\QQ{3} نعتبر في المستوي المركب المنسوب إلى المعلم المتعامد والمتجانس $\left(O;\vec u,\vec v\right)$ النقط $A$، $B$، $C$ و$E$ التي لاحقاتها:
$z_A=-2+i$، $z_B=-2-i$، $z_C=\alpha$ و$z_E=\sqrt3$ حيث $\alpha$ عدد حقيقي و$\alpha>-2$.
بين أن قيمة $\alpha$ التي من أجلها يكون المثلث $ABC$ متقايس الأضلاع هي $\left(-2+\sqrt3\right)$.
نضع في كل ما يأتي: $z_C=-2+\sqrt3$.

\QQ{4} أكتب العدد المركب $\dfrac{z_C-z_E}{z_A-z_B}$ على الشكل الأسي ثم استنتج أن:
\begin{sq}
\qa المستقيمين $(AB)$ و$(EC)$ متعامدان.
\qb النقط $A$، $B$ و$E$ تنتمي إلى نفس الدائرة $(\gamma)$ وعين مركزها ونصف قطرها.
\end{sq}
\QQ{5} ليكن $r$ الدوران الذي يحول النقطة $B$ إلى $C$ ويحول $C$ إلى $A$، عبارته المركبة هي:
$z'=az+\left(\dfrac{\sqrt3-6}{2}\right)+i\left(\dfrac{2\sqrt3-1}{2}\right)$.
\begin{sq}
\qa أحسب العدد المركب $a$، ثم استنتج زاوية الدوران $r$.
\qb تحقق أن النقطة $G$ مركز ثقل المثلث $ABC$ هي مركز الدوران $r$.
\end{sq}
\end{tamrin}

% ===================== التمرين 12 =====================
\begin{tamrin}{بكالوريا 2017 (الدورة الاستدراكية) -- الموضوع الأول -- شعبة رياضيات}
\textbf{I)} حل في مجموعة الأعداد المركبة $\mathbb{C}$ المعادلة ذات المجهول $z$: $2z^2-10z+\dfrac{29}{2}=0$.

\textbf{II)} المستوي المركب منسوب إلى المعلم المتعامد والمتجانس $\left(O;\vec u,\vec v\right)$.
نعتبر النقط $A$، $B$، $C$ و$D$ التي لاحقاتها: $z_A=\dfrac32+\sqrt2\,e^{i\frac{\pi}{4}}$، $z_B=\dfrac32\,e^{-i\frac{\pi}{2}}$، $z_C=-\overline{z_A}$ و$z_D=i$.

\QQ{1}
\begin{sq}
\qa أكتب كلا من $z_A$، $z_B$ و$z_C$ على الشكل الجبري، ثم علم $A$، $B$، $C$ و$D$ في المعلم السابق.
\qb أكتب العدد المركب $\dfrac{z_A-z_B}{z_C-z_B}$ على الشكل الأسي ثم استنتج طبيعة المثلث $ABC$.
\end{sq}
\QQ{2} جد لاحقة النقطة $E$ نظيرة $B$ بالنسبة للنقطة $D$، ثم استنتج طبيعة الرباعي $ABCE$.

\QQ{3} أكتب العبارة المركبة للتشابه $S$ الذي مركزه $B$ ويحول $A$ إلى $D$ ثم حدد نسبته وزاويته.

\end{tamrin}

% ===================== التمرين 13 =====================
\begin{tamrin}{بكالوريا 2017 (الدورة الاستدراكية) -- الموضوع الثاني -- شعبة رياضيات}
\textbf{I)} نعتبر في مجموعة الأعداد المركبة $\mathbb{C}$ المعادلة $(E)$: $z^2-2(1-\sin\alpha)z+2(1-\sin\alpha)=0$، حيث $\alpha$ عدد حقيقي.
(نرمز بـ $z_1$ و$z_2$ إلى حلي المعادلة $(E)$.)

\QQ{1} عين الحلين $z_1$ و$z_2$ بدلالة $\alpha$.

\QQ{2} نضع: $\alpha=\dfrac{\pi}{6}$. بين أن: $z_1^{\,2017}+z_2^{\,2017}=1$.

\textbf{II)} المستوي المركب منسوب إلى المعلم المتعامد والمتجانس $\left(O;\vec u,\vec v\right)$.
نعتبر النقط $A$، $B$ و$C$ التي لاحقاتها: $z_A=\dfrac12+\dfrac{\sqrt3}{2}i$، $z_B=\overline{z_A}$ و$z_C=2z_A$.

\QQ{1} عين قيم العدد الطبيعي $n$ التي يكون من أجلها $\left(\dfrac{z_A}{z_B}\right)^n$ عددا حقيقيا موجبا.

\QQ{2} ليكن $S$ التحويل النقطي الذي يرفق بكل نقطة $M$ ذات اللاحقة $z$ النقطة $M'$ ذات اللاحقة $z'$ والمعرف بـ: $z'=\left(1+z_A\right)z+2z_B$.
عين طبيعة التحويل $S$، ثم حدد عناصره المميزة.

\QQ{3} $(\Gamma)$ مجموعة النقط $M$ ذات اللاحقة $z$ حيث: $\arg\left(\overline{z}-z_B\right)=-\dfrac{\pi}{3}+2k\pi$ و$k\in\mathbb{Z}$.
تحقق أن النقطة $C$ تنتمي إلى $(\Gamma)$، ثم حدد طبيعة المجموعة $(\Gamma)$ وأنشئها.
\end{tamrin}

% ===================== التمرين 14 =====================
\begin{tamrin}{بكالوريا 2017 -- الموضوع الأول -- شعبة رياضيات}
\QQ{1} حل في مجموعة الأعداد المركبة $\mathbb{C}$ المعادلة ذات المجهول $z$: $\left(z-2+2i\right)\left(z^2-2\sqrt2z+8\right)=0$.

\QQ{2} المستوي المركب منسوب إلى المعلم المتعامد والمتجانس $\left(O;\vec u,\vec v\right)$.
نعتبر النقط $A$، $B$ و$C$ التي لاحقاتها: $z_A=\sqrt2+i\sqrt6$، $z_B=\overline{z_A}$ و$z_C=2(1-i)$.
\begin{sq}
\qa أكتب كلا من $z_A$، $z_B$ و$z_C$ على الشكل الأسي، ثم استنتج أن النقط $A$، $B$ و$C$ تنتمي إلى دائرة $(\Omega)$ يطلب تعيين مركزها ونصف قطرها.
\qb عين قيم العدد الطبيعي $n$ التي من أجلها يكون العدد المركب $\left(\dfrac{z_A}{z_C}\right)^n$ تخيليا صرفا.
\qc نسمي $(\Gamma)$ مجموعة النقط $M$ من المستوي ذات اللاحقة $z$ حيث $z=z_C-k\left(\dfrac{z_A}{z_B}\right)$ مع $k$ يمسح $\mathbb{R}_+$.
تحقق أن النقطة $C$ تنتمي إلى $(\Gamma)$، ثم عين وأنشئ $(\Gamma)$.
\end{sq}
\QQ{3} $r$ الدوران الذي مركزه $O$ وزاويته $\dfrac{2\pi}{3}$ و$h$ التحاكي الذي مركزه $O$ ونسبته $-2$.
عين طبيعة التحويل $h\circ r$ وعناصره المميزة، ثم استنتج صورة الدائرة $(\Omega)$ بالتحويل $h\circ r$.
\end{tamrin}

% ===================== التمرين 15 =====================
\begin{tamrin}{بكالوريا 2017 -- الموضوع الثاني -- شعبة رياضيات}
\textbf{I)} أكتب العدد $\left(\dfrac52+i\right)^2$ على الشكل الجبري.

\textbf{II)} المستوي المركب منسوب إلى المعلم المتعامد والمتجانس $\left(O;\vec u,\vec v\right)$.
نعتبر النقط $A$، $B$، $C$ و$I$ ذات اللواحق: $z_A=\dfrac32+\sqrt2\,e^{i\frac{\pi}{4}}$، $z_B=-\dfrac32i$، $z_C=-\overline{z_A}$ و$z_I=i$.

\QQ{1} أكتب كلا من $z_A$ و$z_C$ على الشكل الجبري.

\QQ{2} أكتب العدد المركب $\dfrac{z_C-z_B}{z_A-z_B}$ على الشكل الأسي مستنتجا طبيعة المثلث $ABC$.

\QQ{3} ليكن $S$ التشابه المباشر الذي مركزه $B$ ويحول $A$ إلى $I$.
\begin{sq}
\qa أكتب العبارة المركبة للتشابه $S$، ثم عين نسبته وزاويته.
\qb نعرف من أجل كل عدد طبيعي $n$ حيث $n\geq2$، التحويل النقطي $T_n$ كما يلي: $T_n=\underbrace{S\circ S\circ\cdots\circ S}_{n\ \text{fois}}$.
عين قيم العدد الطبيعي $n$ حتى يكون $T_n$ تحاكيا، عين عندئذ عناصره المميزة.
\end{sq}
\end{tamrin}

% ===================== التمرين 16 =====================
\begin{tamrin}{بكالوريا 2016 -- الموضوع الأول -- شعبة رياضيات}
\QQ{1}
\begin{sq}
\qa حل في $\mathbb{C}$ المعادلة: $z^2-4z+5=0$.
\qb استنتج حلول المعادلة ذات المجهول $z$ الآتية: $\left(z+1+i\left(1-\sqrt3\right)\right)^2-4z+1-4i\left(1-\sqrt3\right)=0$.
\end{sq}
\QQ{2} $\theta$ عدد حقيقي حيث $0\leq\theta\leq\dfrac{\pi}{2}$، و$z_0$ عدد مركب طويلته $1$ و$\theta$ عمدة له.
\begin{sq}
\qa أكتب العدد المركب $1+i\sqrt3$ على الشكل الأسي.
\qb عين $\theta$ علما أن: $\dfrac{z_0\left(1+i\sqrt3\right)}{\overline{z_0}}=2e^{i\frac{\pi}{2}}$ ($\overline{z_0}$ مرافق العدد المركب $z_0$).
\qc $n$ عدد طبيعي. من أجل قيمة $\theta$ المتحصل عليها أكتب العدد المركب $\left[\dfrac{z_0\left(1+i\sqrt3\right)}{2}\right]^n$ على الشكل المثلثي.
\qd عين قيم العدد الطبيعي $n$ التي من أجلها يكون $\left[\dfrac{z_0\left(1+i\sqrt3\right)}{2}\right]^n$ عددا حقيقيا موجبا تماما.
\end{sq}
\QQ{3} المستوي المركب منسوب إلى المعلم المتعامد والمتجانس $\left(O;\vec u,\vec v\right)$.
نعتبر النقط $A$، $B$ و$C$ التي لاحقاتها على الترتيب: $z_A=2-i$، $z_B=2+i$ و$z_C=1+i\sqrt3$.
\begin{sq}
\qa عين $z_D$ لاحقة النقطة $D$ مرجح الجملة $\left\{(A,1);(B,-1);(C,1)\right\}$.
\qb استنتج أن الرباعي $ABCD$ متوازي أضلاع.
\qc $E$ النقطة من المستوي المركب ذات اللاحقة $z_E$ حيث:
$\begin{cases}\arg\left(z_E-z_A\right)-\arg\left(z_E-z_B\right)=\dfrac{\pi}{2}\\[4pt] \left|\dfrac{z_E-z_A}{z_E-z_B}\right|=2\end{cases}$

-- بين أن: $z_E=\dfrac{14}{5}+\dfrac35i$.\par
-- بين أن النقطة $A$ هي صورة النقطة $B$ بتشابه مباشر يطلب تعيين عناصره المميزة.
\end{sq}
\QQ{4} $M$ نقطة من المستوي المركب لاحقتها $z$، $I$ منتصف القطعة المستقيمة $[AB]$.
\begin{sq}
\qa عين $z_I$ لاحقة النقطة $I$.
\qb $\alpha$ عدد حقيقي. نسمي $(\Gamma)$ مجموعة النقط $M$ من المستوي التي تحقق: $z-z_I=e^{i\alpha}$.

-- تحقق أن النقطة $E$ تنتمي للمجموعة $(\Gamma)$.\par
-- عين طبيعة المجموعة $(\Gamma)$ وعناصرها المميزة عندما يتغير $\alpha$ في $\mathbb{R}$.
\end{sq}
\end{tamrin}

% ===================== التمرين 17 =====================
\begin{tamrin}{بكالوريا 2016 -- الموضوع الثاني -- شعبة رياضيات}
\textbf{I)}\QQ{1} حل في مجموعة الأعداد المركبة $\mathbb{C}$ المعادلة: $z^2-2z+2=0$.

\QQ{2} جد العددين المركبين $z_1$ و$z_2$ حيث: $\begin{cases}2z_1-3z_2=5i\sqrt2\\ z_1+3z_2=-2i\sqrt2\end{cases}$

\textbf{II)} المستوي المركب منسوب إلى المعلم المتعامد والمتجانس $\left(O;\vec u,\vec v\right)$.
النقط $A$، $B$، $C$، $D$ و$H$ لاحقاتها على الترتيب: $z_A=i\sqrt2$، $z_B=-i\sqrt2$، $z_C=1+i$، $z_D=1-i$ و$z_H=\dfrac{z_C-z_B}{z_E-z_B}$،
حيث $E$ النقطة التي تحقق: $\overrightarrow{DE}=2\overrightarrow{DO}$.

\QQ{1} أكتب $z_H$ على الشكل الأسي واستنتج نوع المثلث $BEC$.

\QQ{2} $S$ تحويل نقطي في المستوي يرفق بكل نقطة $M$ لاحقتها $z$ النقطة $M'$ لاحقتها $z'$ حيث: $z'=z_A\cdot z+z_B$.
\begin{sq}
\qa ما هي طبيعة التحويل $S$؟ وما هي عناصره المميزة؟
\qb أحسب مساحة الدائرة $(\gamma)$ التي مركزها $C$ ونصف قطرها $CD$.
\qc عين $(\gamma')$ صورة $(\gamma)$ بالتحويل $S$ واستنتج مساحتها.
\end{sq}
\QQ{3} عين $(\delta)$ مجموعة النقط $M$ من المستوي ($M$ تختلف عن $B$ و$C$) ذات اللاحقة $z$ التي يكون من أجلها العدد $\dfrac{z_B-z}{z_C-z}$ حقيقيا سالبا.
\end{tamrin}

% ===================== التمرين 18 =====================
\begin{tamrin}{بكالوريا 2015 -- الموضوع الأول -- شعبة رياضيات}
ينسب المستوي إلى المعلم المتعامد والمتجانس $\left(O;\vec u,\vec v\right)$. نعتبر النقط $A$، $B$، $C$، $H$ و$I$ لاحقاتها على الترتيب:
$z_A=i$، $z_B=-2+i$، $z_C=-3$، $z_H=-3+4i$ و$z_I=-1-i$.

\QQ{1}
\begin{sq}
\qa أعلم النقط $A$، $B$، $C$، $H$ و$I$.
\qb عين النسبة والزاوية للتشابه المباشر الذي مركزه $B$ ويحول النقطة $A$ إلى النقطة $C$.
\end{sq}
\QQ{2} عين $z_G$ لاحقة النقطة $G$ مركز ثقل المثلث $ABC$.

\QQ{3}
\begin{sq}
\qa أكتب على الشكل الجبري العدد المركب $\dfrac{z_B-z_C}{z_H-z_A}$.
\qb استنتج أن المستقيمين $(AH)$ و$(BC)$ متعامدان.
\qc بين أن $H$ هي نقطة تلاقي ارتفاعات المثلث $ABC$.
\end{sq}
\QQ{4} بين أن النقط $H$، $G$ و$I$ في استقامية.

\QQ{5} $(\Gamma)$ مجموعة النقط $M$ من المستوي ذات اللاحقة $z$ حيث: $z+1+i=\sqrt5\,e^{i\theta}$ مع $\theta\in\mathbb{R}$.
\begin{sq}
\qa بين أن النقطة $A$ تنتمي إلى المجموعة $(\Gamma)$.
\qb عين طبيعة المجموعة $(\Gamma)$ مع تحديد عناصرها المميزة.
\qc أنشئ المجموعة $(\Gamma)$.
\qd تحقق أن النقطتين $B$ و$C$ تنتميان إلى المجموعة $(\Gamma)$.
\end{sq}
\end{tamrin}

% ===================== التمرين 19 =====================
\begin{tamrin}{بكالوريا 2015 -- الموضوع الثاني -- شعبة رياضيات}
\QQ{1} حل في مجموعة الأعداد المركبة $\mathbb{C}$ المعادلة ذات المجهول $z$ التالية: $z^2-2\left(1-\sqrt3\right)z+8=0$.
(لاحظ أن: $\left(1+\sqrt3\right)^2=4+2\sqrt3$.)

\QQ{2} المستوي المركب منسوب إلى المعلم المتعامد والمتجانس $\left(O;\vec u,\vec v\right)$.
$A$ و$B$ نقطتان من المستوي، لاحقتاهما على الترتيب: $z_A=\left(1-\sqrt3\right)+i\left(1+\sqrt3\right)$ و$z_B=\overline{z_A}$.
\begin{sq}
\qa بين أن: $\dfrac{z_B}{z_A}=e^{-i\frac{7\pi}{6}}$.
\qb استنتج عمدة للعدد المركب $z_A$.
\qc استنتج القيمة المضبوطة لكل من العددين $\cos\dfrac{7\pi}{12}$ و$\sin\dfrac{7\pi}{12}$.
\end{sq}
\QQ{3}
\begin{sq}
\qa حل، في مجموعة الأعداد الصحيحة، المعادلة ذات المجهول $(x;y)$ التالية: $7x-2y=1$.
\qb بين أنه إذا كانت الثنائية $(x;y)$ من الأعداد الصحيحة حلا للمعادلة $7x-24y=12$ فإن $x$ يكون مضاعفا للعدد $12$.
\qc استنتج كل الثنائيات $(x;y)$ من الأعداد الصحيحة حلول للمعادلة $7x-24y=12$.
\qd عين مجموعة قيم العدد الطبيعي $n$ التي يكون من أجلها العدد $\left(z_A\right)^n$ عددا حقيقيا سالبا تماما.
\end{sq}
\end{tamrin}
% ---------- helpers for figures (part m2) ----------
\providecommand{\MtwoAx}[4]{%
 \draw[gray!20,very thin] (#1,#3) grid (#2,#4);
 \draw[-{Stealth[length=2mm]}] (#1,0)--(#2,0) node[below left]{$x$};
 \draw[-{Stealth[length=2mm]}] (0,#3)--(0,#4) node[below left]{$y$};
 \node[below left,font=\scriptsize] at (0,0){$O$};
 \draw[-{Stealth[length=1.6mm]},thick] (0,0)--(1,0) node[below,font=\scriptsize]{$\vec u$};
 \draw[-{Stealth[length=1.6mm]},thick] (0,0)--(0,1) node[left,font=\scriptsize]{$\vec v$};}
\providecommand{\MtwoPt}[3]{\fill (#1) circle (1.6pt) node[#2,font=\small]{$#3$};}

% ======================= التمرين 20 =======================
\begin{tamrin}{بكالوريا 2014 -- الموضوع الأول -- شعبة رياضيات}
\QQ{1} حل في مجموعة الأعداد المركبة $\mathbb{C}$ المعادلة التالية:
\[(z-1-2i)\left(z^2-2\left(1+\sqrt3\right)z+5+2\sqrt3\right)=0.\]
\QQ{2} $A$، $B$، $C$ و$D$ نقط من المستوي المركب المنسوب إلى المعلم المتعامد والمتجانس $(O;\vec u,\vec v)$ لاحقاتها على الترتيب:
\[z_A=1+2i,\qquad z_B=1+\sqrt3+i,\qquad z_C=1+\sqrt3-i,\qquad z_D=1-2i.\]
\begin{sq}
\qa بيّن أن: $AB=CD$ و$(AD)$ يوازي $(BC)$.
\qb تحقق أن: $\dfrac{z_B+z_D}{2}\neq\dfrac{z_A+z_C}{2}$ واستنتج طبيعة الرباعي $ABCD$.
\end{sq}
\QQ{3}
\begin{sq}
\qa بيّن أن: $\dfrac{z_D-z_B}{z_A-z_B}=\sqrt3\,e^{i\frac{\pi}{2}}$، استنتج أن $D$ صورة $A$ بتشابه مباشر مركزه $B$ يطلب تعيين نسبته وزاويته.
\qb بيّن أن المثلث $ADB$ قائم وأن النقط $A$، $B$، $C$ و$D$ تنتمي لدائرة يطلب تحديد مركزها ونصف قطرها.
\qc استنتج إنشاء للرباعي $ABCD$.
\end{sq}
\end{tamrin}

% ======================= التمرين 21 =======================
\begin{tamrin}{بكالوريا 2014 -- الموضوع الثاني -- شعبة رياضيات}
المستوي منسوب إلى المعلم المتعامد والمتجانس $(O;\vec u,\vec v)$. $A$ و$B$ النقطتان اللتان لاحقتاهما على الترتيب: $a=-2+6i$ و$b=-1+2i$.
\QQ{1} أكتب العدد المركب $1+i$ على الشكل الأسي.
\QQ{2} $S$ التحويل النقطي الذي يرفق بكل نقطة $M$ لاحقتها $z$ النقطة $M'$ لاحقتها $z'$ حيث: $z'=\sqrt2\,e^{i\frac{\pi}{4}}z+2$.
\begin{sq}
\qa $D$ النقطة ذات اللاحقة $d$ حيث $d=2i$. جد لاحقة النقطة $D'$ صورة $D$ بالتحويل $S$، ماذا تستنتج؟
\qb بيّن أن: $z'-d=\sqrt2\,e^{i\frac{\pi}{4}}(z-d)$ واستنتج طبيعة وعناصر التحويل $S$.
\end{sq}
\QQ{3} $(\Delta)$ المستقيم ذو المعادلة: $3x+5y=11$.
\begin{sq}
\qa تحقق أن النقطة $M_0(-3;4)$ تنتمي إلى $(\Delta)$، ثم عيّن نقط $(\Delta)$ التي إحداثياتها أعداد صحيحة.
\qb $M_0'$ صورة $M_0$ بالتحويل $S$، بيّن أن المستقيمين $(BM_0')$ و$(BA)$ متعامدان.
\end{sq}
\QQ{4} $x$ و$y$ عددان صحيحان من المجال $[-5;5]$. عيّن مجموعة النقط $M(x;y)$ من المستوي بحيث يكون المستقيمان $(BA)$ و$(BM')$ متعامدين، حيث $M'$ هي صورة $M$ بالتحويل $S$.
\end{tamrin}

% ======================= التمرين 22 =======================
\begin{tamrin}{بكالوريا 2013 -- الموضوع الأول -- شعبة رياضيات}
\textbf{I)} $a$ و$b$ عددان موجبان تمامًا. نعتبر في المستوي المنسوب إلى المعلم المتعامد والمتجانس $(O;\vec u,\vec v)$ النقط $A$، $B$، $C$ و$E$ التي لاحقاتها على الترتيب:
\[z_A=a\,e^{i\frac{3\pi}{4}},\qquad z_B=-a\sqrt2,\qquad z_C=\overline{z_A},\qquad z_E=b\,e^{i\frac{3\pi}{2}}.\]
\QQ{1}
\begin{sq}
\qa أكتب على الشكل الأسي العدد المركب $\dfrac{z_A-z_B}{z_A}$، ثم استنتج طبيعة المثلث $OAB$.
\qb حدد طبيعة الرباعي $OABC$، ثم استنتج مساحته.
\end{sq}
\QQ{2} التشابه المباشر $S$ ذو المركز $O$ والنسبة $\dfrac{b}{a}$ والزاوية $\dfrac{3\pi}{4}$، يحول كل نقطة $M(z)$ إلى النقطة $M'(z')$.
\begin{sq}
\qa أكتب العبارة المركبة للتشابه المباشر $S$، ثم تحقق أن: $S(A)=E$.
\qb بيّن أن مساحة الرباعي $OEFG$ هي $b^2$ (مقدرة بوحدة المساحة) حيث: $S(B)=F$ و$S(C)=G$.
\end{sq}
\QQ{3}
\begin{sq}
\qa أحسب بدلالة $a$ و$b$ العبارة: $|z_C|^2+|z_E|^2-2\,|z_C\times z_E|\cos\left[\arg\left(\dfrac{z_C}{z_E}\right)\right]$.
\qb استنتج $CE^2$ بدلالة $a$ و$b$.
\end{sq}
\end{tamrin}

% ======================= التمرين 23 =======================
\begin{tamrin}{بكالوريا 2013 -- الموضوع الثاني -- شعبة رياضيات}
\QQ{1} حل في مجموعة الأعداد المركبة $\mathbb C$ المعادلة التالية: $z^2+z+1=0$.
\QQ{2} نعتبر في المستوي المركب المنسوب إلى المعلم المتعامد والمتجانس $(O;\vec u,\vec v)$، النقط $A$، $B$ و$M$ ذات اللاحقات $z_A=-\dfrac{1+i\sqrt3}{2}$، $z_B=\overline{z_A}$ و$z$ على الترتيب.
\begin{sq}
\qa أكتب $z_A$ على الشكل الأسي.
\qb عيّن مجموعة النقط $M$ من المستوي، حيث: $\arg\left[(z-z_A)^2\right]=\arg(z_A)-\arg(z_B)$.
\end{sq}
\QQ{3}
\begin{sq}
\qa التحويل النقطي $r$، يرفق بكل نقطة $M(z)$ النقطة $M'(z')$، حيث: $z'=z_A\,z+z_B\sqrt3$. ما طبيعة التحويل $r$؟ عيّن عناصره المميزة.
\qb عيّن نسبة ومركز التحاكي $h$، يرفق بكل نقطة $M(z)$ إلى النقطة $M'(z')$، حيث: $z'=-2z+3i$.
\qc نضع: $S=h\circ r$. عيّن طبيعة التحويل $S$ مبرزًا عناصره المميزة، ثم تحقق أن عبارته المركبة هي: $z'=2e^{i\frac{\pi}{3}}(z-i)+i$.
\end{sq}
\QQ{4} نعتبر النقطة $\Omega$ ذات اللاحقة $i$ والنقط $C$، $D$ و$E$ حيث: $S(O)=C$، $S(C)=D$ و$S(D)=E$.
بيّن أن النقط $O$، $\Omega$ و$E$ في استقامية.
\QQ{5}
\begin{sq}
\qa عيّن $(\Gamma)$ مجموعة النقط $M(z)$ حيث $z=2e^{i\theta}+e^{i\frac{\pi}{2}}$ مع $\theta\in\R$.
\qb عيّن $(\Gamma')$ صورة $(\Gamma)$ بالتحويل $S$.
\end{sq}
\end{tamrin}

% ======================= التمرين 24 =======================
\begin{tamrin}{بكالوريا 2012 -- الموضوع الأول -- شعبة رياضيات}
\QQ{1} حل في مجموعة الأعداد المركبة $\mathbb C$ المعادلة ذات المجهول $z$ التالية: $z^2-\sqrt2\,z+1=0$.
\QQ{2} المستوي المركب منسوب إلى المعلم المتعامد والمتجانس $(O;\vec u,\vec v)$. $A$، $B$، $C$ نقط من المستوي لاحقاتها على الترتيب: $z_A=\dfrac{1+i}{\sqrt2}$، $z_B=\overline{z_A}$ و$z_C=z_A+z_B$.
\begin{sq}
\qa أكتب على الشكل الأسي الأعداد المركبة: $z_A$، $z_B$ و$\dfrac{z_A}{z_B}$.
\qb عيّن لاحقة $A'$، $B'$ و$C'$ صور النقط $A$، $B$ و$C$ على التوالي بالدوران الذي مركزه $O$ وزاويته $\dfrac{\pi}{4}$.
\qc بيّن أن الرباعي $OA'C'B'$ مربع.
\end{sq}
\QQ{3} نسمي $(\Delta)$ مجموعة النقط $M$ من المستوي ذات اللاحقة $z$ حيث: $|z-z_A|=|z-z_B|$.
\begin{sq}
\qa بيّن أن $(\Delta)$ هو محور الفواصل.
\qb بيّن أن حلي المعادلة: $\left(\dfrac{z-z_A}{z-z_B}\right)^2=i$ عددان حقيقيان. (لا يُطلب حساب الحلين)
\end{sq}
\end{tamrin}

% ======================= التمرين 25 =======================
\begin{tamrin}{بكالوريا 2012 -- الموضوع الثاني -- شعبة رياضيات}
\QQ{1} حل في مجموعة الأعداد المركبة $\mathbb C$ المعادلة ذات المجهول $z$ التالية: $(z^2+4)(z^2-2\sqrt3\,z+4)=0$.
\QQ{2} نعتبر في المستوي المركب المنسوب إلى المعلم المتعامد والمتجانس $(O;\vec u,\vec v)$ النقط $A$، $B$، $C$ و$D$ التي لواحقها على الترتيب:
$z_A=\sqrt3+i$، $z_B=\overline{z_A}$، $z_C=-2i$ و$z_D=\overline{z_C}$.
بيّن أن النقط $A$، $B$، $C$ و$D$ تنتمي إلى دائرة $(\gamma)$ يطلب تعيين مركزها ونصف قطرها، ثم أنشئ النقط $A$، $B$، $C$ و$D$.
\QQ{3} نرمز بـ $z_E$ إلى لاحقة النقطة $E$ نظيرة النقطة $B$ بالنسبة إلى المبدأ $O$.
\begin{sq}
\qa بيّن أن: $\dfrac{z_A-z_C}{z_E-z_C}=e^{-i\frac{\pi}{3}}$.
\qb بيّن أن النقطة $A$ هي صورة النقطة $E$ بدوران $R$ مركزه $C$ يطلب تعيين زاويته.
\qc استنتج طبيعة المثلث $AEC$.
\qd $H$ هو التحاكي الذي مركزه $O$ ونسبته $2$. عيّن طبيعة التحويل $R\circ H$ وعناصره المميزة، ثم استنتج صورة الدائرة $(\gamma)$ بالتحويل $R\circ H$.
\end{sq}
\end{tamrin}

% ======================= التمرين 26 =======================
\begin{tamrin}{بكالوريا 2011 -- الموضوع الأول -- شعبة رياضيات}
المستوي منسوب إلى المعلم المتعامد والمتجانس $(O;\vec u,\vec v)$. $A$، $B$، $C$ ثلاث نقط من المستوي لاحقاتها على الترتيب: $z_A=1-i$، $z_B=-1+i$، $z_C=\sqrt3(1+i)$.
\QQ{1} أكتب على الشكل الأسي الأعداد المركبة: $z_A$، $z_B$ و$z_C$.
\QQ{2}
\begin{sq}
\qa أحسب الطويلة وعمدة للعدد المركب $\dfrac{z_B-z_A}{z_C-z_A}$، ثم فسر النتائج المحصل عليها.
\qb حدد طبيعة المثلث $ABC$.
\end{sq}
\QQ{3} عيّن لاحقة النقطة $D$ بحيث يكون الرباعي $ACBD$ معينًا.
\QQ{4} $T$ التحويل النقطي الذي يرفق بكل نقطة $M$ من المستوي لاحقتها $z$ النقطة $M'$ ذات اللاحقة $z'$ حيث: $z'=(-1+i)z+1-3i$.
\begin{sq}
\qa عيّن طبيعة التحويل $T$ وعناصره المميزة.
\qb استنتج طبيعة التحويل $T\circ T$ وعناصره المميزة.
\end{sq}
\end{tamrin}

% ======================= التمرين 27 =======================
\begin{tamrin}{بكالوريا 2011 -- الموضوع الثاني -- شعبة رياضيات}
أجب بصحيح أو خطأ مع التبرير في كل حالة من الحالات الآتية:
\QQ{1}
\begin{sq}
\qa الشكل المثلثي للعدد المركب $a=-\dfrac{\sqrt2}{2}+i\dfrac{\sqrt2}{2}$ هو $-\cos\dfrac{\pi}{4}+i\sin\dfrac{\pi}{4}$.
\qb $a^{2011}+\overline a=0$، حيث $\overline a$ مرافق $a$.
\end{sq}
\QQ{2} في المستوي المركب المنسوب إلى المعلم المتعامد والمتجانس $(O;\vec u,\vec v)$.
\begin{sq}
\qa التحويل الذي كتابته المركبة: $z'=\left(-\dfrac{\sqrt2}{2}+i\dfrac{\sqrt2}{2}\right)z$ دوران زاويته $-\dfrac{\pi}{4}$ ومركزه مبدأ المعلم.
\qb مجموعة النقط $M$ ذات اللاحقة $z$ حيث: $\arg(z-i)=-\dfrac{\pi}{4}$ هي المستقيم $(\Delta)$ الذي يشمل النقطة $A$ ذات اللاحقة $i$ وشعاع توجيهه $\vec w$ لاحقته $1+i$.
\end{sq}
\end{tamrin}
\begin{molahada}
في نص الموضوع رُمز لشعاع التوجيه بـ $\vec u$، وهو نفس رمز شعاع المعلم؛ سميناه $\vec w$ لرفع الالتباس.
\end{molahada}

% ======================= التمرين 28 =======================
\begin{tamrin}{بكالوريا 2010 -- الموضوع الأول -- شعبة رياضيات}
نعتبر في مجموعة الأعداد المركبة $\mathbb C$ المعادلة: $(E)\ldots\ z^3-3z^2+3z-9=0$.
\QQ{1}
\begin{sq}
\qa تحقق أن $3$ حل للمعادلة $(E)$، ثم عيّن الأعداد الحقيقية $a$، $b$ و$c$ بحيث يكون من أجل كل عدد مركب $z$ فإن: $z^3-3z^2+3z-9=(z-3)(az^2+bz+c)$.
\qb حل في $\mathbb C$ المعادلة $(E)$.
\end{sq}
\QQ{2} المستوي منسوب إلى المعلم المتعامد والمتجانس $(O;\vec u,\vec v)$. النقط $A$، $B$ و$C$ صور الأعداد $z_A=3$، $z_B=\sqrt3\,i$، $z_C=-\sqrt3\,i$. بيّن أن المثلث $ABC$ متقايس الأضلاع.
\QQ{3} $D$ النقطة التي لاحقتها $z_D=2e^{i\frac{5\pi}{6}}$ و$E$ صورتها بالدوران الذي مركزه $O$ وزاويته $\dfrac{\pi}{3}$. عيّن $z_E$ لاحقة النقطة $E$.
\QQ{4} $F$ النقطة التي لاحقتها $z_F=1-\sqrt3\,i$.
\QQ{5}
\begin{sq}
\qa أحسب $\dfrac{z_F}{z_E}$ واستنتج أن المستقيمين $(OE)$ و$(OF)$ متعامدان.
\qb عيّن $z_G$ لاحقة النقطة $G$ بحيث يكون $OEGF$ مربعًا.
\end{sq}
\end{tamrin}

% ======================= التمرين 29 =======================
\begin{tamrin}{بكالوريا 2010 -- الموضوع الثاني -- شعبة رياضيات}
المستوي منسوب إلى المعلم المتعامد والمتجانس $(O;\vec u,\vec v)$.
\QQ{1} نسمي $A$، $B$، $I$ النقط التي لاحقاتها على الترتيب: $z_A=1-4i$، $z_B=-1-2i$، $z_I=1-2i$.
\begin{sq}
\qa علّم النقط $A$، $B$، $I$.
\qb أكتب على الشكل الجبري العدد المركب $z=\dfrac{z_I-z_A}{z_I-z_B}$.
\qc ما نوع المثلث $IAB$؟
\qd $C$ صورة $I$ بالتحاكي الذي مركزه $A$ ونسبته $2$. أحسب اللاحقة $z_C$ للنقطة $C$.
\qe $D$ مرجح $\{(A,1);(B,-1);(C,1)\}$. أحسب اللاحقة $z_D$ للنقطة $D$.
\item[(و)] بيّن أن $ABCD$ مربع.
\end{sq}
\QQ{2} عيّن وأنشئ $(\Gamma_1)$ مجموعة النقط $M$ من المستوي حيث: $\left\|\overrightarrow{MA}-\overrightarrow{MB}+\overrightarrow{MC}\right\|=\dfrac12\left\|\overrightarrow{MA}+\overrightarrow{MC}\right\|$.
\QQ{3} عيّن وأنشئ $(\Gamma_2)$ مجموعة النقط $M$ من المستوي حيث: $\left\|\overrightarrow{MA}-\overrightarrow{MB}+\overrightarrow{MC}\right\|=1$.
\end{tamrin}

% ======================= التمرين 30 =======================
\begin{tamrin}{بكالوريا 2009 -- الموضوع الثاني -- شعبة رياضيات}
نرفق بكل عدد مركب $z$ يختلف عن $1$ العدد المركب $f(z)$ حيث: $f(z)=\dfrac{z-i}{z-1}$.
\QQ{1} حل في مجموعة الأعداد المركبة $\mathbb C$ المعادلة: $(45+45i)f(z)=23+45i-2z$.
\QQ{2} لتكن $M$ صورة العدد المركب $z$ في المستوي المنسوب إلى المعلم المتعامد والمتجانس $(O;\vec u,\vec v)$.
\begin{sq}
\qa عيّن مجموعة النقط $M$ من المستوي بحيث يكون $f(z)$ حقيقيًا سالبًا تمامًا.
\qb أحسب العدد المركب $z_0$ بحيث: $|f(z_0)|=1$ و$\arg\left(f(z_0)\right)=\dfrac{3\pi}{2}$.
\end{sq}
\QQ{3} في المستوي المركب نعتبر النقط $A$، $B$ و$C$ صور الأعداد المركبة $1$، $i$ و$z_0$ على الترتيب.
\begin{sq}
\qa ما نوع المثلث $ABC$؟
\qb عيّن النقطة $D$ نظيرة $C$ بالنسبة إلى المستقيم $(AB)$ واستنتج طبيعة الرباعي $ACBD$.
\end{sq}
\end{tamrin}

% ======================= التمرين 31 =======================
\begin{tamrin}{بكالوريا 2008 -- الموضوع الأول -- شعبة رياضيات}
المستوي منسوب إلى المعلم المتعامد والمتجانس $(O;\vec u,\vec v)$. نعتبر النقطتين $A$ و$B$ اللتين لاحقتاهما $\sqrt3-i$ و$\sqrt3+3i$ على الترتيب.
\QQ{1} أكتب العبارة المركبة للتشابه المباشر $S$ الذي مركزه $O$ ويحول $A$ إلى $B$، ثم عيّن زاويته ونسبته.
\end{tamrin}

% ======================= التمرين 32 =======================
\begin{tamrin}{بكالوريا 2008 -- الموضوع الثاني -- شعبة رياضيات}
نعتبر في مجموعة الأعداد المركبة $\mathbb C$ كثير الحدود $P(z)$ المعرف كما يلي: $P(z)=2z^4-2iz^3-z^2-2iz+2$.
\QQ{1} بيّن أنه إذا كان $\alpha$ جذرًا لـ $P(z)$ فإن $\dfrac1\alpha$ جذر له أيضًا.
\QQ{2} تحقق أن: $1+i$ جذر لكثير الحدود $P(z)$.
\QQ{3} في المستوي المركب منسوب إلى المعلم المتعامد والمتجانس $(O;\vec u,\vec v)$ نعتبر النقط $A$، $B$، $C$ و$D$ والتي لواحقها على الترتيب:
$1+i$، $-1+i$، $-\dfrac m2-\dfrac m2i$ و$\dfrac m2-\dfrac m2i$ حيث $m$ عدد حقيقي.
عيّن $m$ حتى يكون الرباعي $ABCD$ مربعًا.
\end{tamrin}
\end{document}
